% Kaelix — revisão da literatura
% Compilar:  latexmk -pdf revisao-literatura.tex   (ou pdflatex duas vezes)

\documentclass[11pt,a4paper]{article}

\usepackage[T1]{fontenc}
\usepackage[utf8]{inputenc}
\usepackage[brazil]{babel}
\usepackage{lmodern}
\usepackage[margin=2.7cm]{geometry}
\usepackage{booktabs}
\usepackage{siunitx}
\usepackage{graphicx}
\usepackage{amsmath}
\usepackage{amssymb}
\usepackage{microtype}
\usepackage[hidelinks]{hyperref}
\usepackage{caption}
\usepackage{enumitem}
\usepackage{xcolor}
\usepackage{titlesec}
\usepackage{float}

% Floats: permitir mais figura por página e evitar páginas de float esvaziadas
\setcounter{topnumber}{2}
\setcounter{bottomnumber}{2}
\setcounter{totalnumber}{3}
\renewcommand{\topfraction}{0.92}
\renewcommand{\bottomfraction}{0.75}
\renewcommand{\textfraction}{0.06}
\renewcommand{\floatpagefraction}{0.80}

% Figuras largas: avançam levemente sobre as margens para preservar legibilidade
\newcommand{\figuralarga}[1]{%
  \makebox[\textwidth][c]{\includegraphics[width=1.13\textwidth]{#1}}}

\sisetup{output-decimal-marker={,}, group-separator={.}, per-mode=symbol}
\captionsetup{font=small, labelfont=bf, width=1.02\textwidth}
\setlist{itemsep=2pt, topsep=4pt, parsep=0pt}
\titleformat{\section}{\normalfont\large\bfseries}{\thesection}{0.7em}{}
\titleformat{\subsection}{\normalfont\normalsize\bfseries}{\thesubsection}{0.7em}{}
\graphicspath{{../../experiments/figures/output/}{../img/}}

\newcommand{\pt}{\,p.p.}
\newcommand{\mgf}{\ensuremath{\mathrm{m}g}}
\newcommand{\dens}{\ensuremath{\mu g/\sqrt{\mathrm{Hz}}}}

\begin{document}

\noindent
\begin{minipage}[c]{0.16\textwidth}
  \includegraphics[width=\linewidth]{skala_logo.jpeg}
\end{minipage}%
\begin{minipage}[c]{0.84\textwidth}
  \centering
  {\LARGE\bfseries Kaelix: revisão da literatura}\\[4pt]
  {\large O que a literatura publicada sustenta, o que ela contesta\\
  e o que ela não cobre nas escolhas de um sensor IoT de vibração de baixo custo}\\[10pt]
  {\normalsize agosto de 2026}
\end{minipage}

\vspace{6pt}
\hrule
\vspace{14pt}

\section{Escopo e método da revisão}

Este documento confronta as escolhas de projeto e os resultados do Kaelix com a
literatura publicada. O Kaelix é um dispositivo IoT de manutenção preditiva para motores
industriais, construído sobre um ESP32-S3-WROOM-1, um acelerômetro MPU6050 em barramento
I\textsuperscript{2}C com taxa de saída de \SI{1}{\kilo\hertz} (e, portanto, frequência de
Nyquist de \SI{500}{\hertz}), um termistor NTC \SI{10}{\kilo\ohm}, rádio LoRa SX1278
em \SI{433}{\mega\hertz}, invólucro em ASA impresso com fixação magnética e bateria LiPo
de \SI{2000}{\milli\ampere\hour}. A detecção de anomalias é feita por Isolation Forest com
inferência nativa em C++ embarcada, sem TFLite Micro, sobre quatro descritores: valor
eficaz, curtose, fator de crista e frequência dominante.

A revisão não é um inventário de trabalhos correlatos. Cada eixo temático parte de um
achado próprio do projeto e pergunta o que a literatura diz sobre ele: se o confirma, se o
contesta ou se não o cobre. Os eixos são:

\begin{enumerate}[label=\textbf{Eixo \arabic*.}, leftmargin=2.6cm]
\item Isolation Forest, limiar de decisão e TinyML embarcado.
\item Protocolo de validação, vazamento entre janelas e confundidores de ensaio.
\item Banda de acelerômetros MEMS e detectabilidade de falha de rolamento.
\item Normalização ISO e integração aceleração--velocidade.
\item Sistemas IoT comparáveis, invólucro e restrições térmicas.
\end{enumerate}

O Eixo~2 não integra esta versão do documento. Nenhuma das referências que o sustentariam
consta do conjunto conferido em fonte primária que ampara este texto, e a regra adotada
aqui foi não citar o que não foi conferido. Os achados próprios correspondentes são dois:
o vazamento por partição aleatória de janela, que infla a acurácia em cerca de 36 pontos
percentuais quando comparado a validação por \emph{GroupKFold} por ensaio, e o confundidor
de massa de desbalanceamento nos ensaios de rolamento do MAFAULDA. Ambos aparecem neste
documento apenas onde incidem sobre os outros quatro eixos, e sempre marcados como
resultado próprio, nunca como literatura. O eixo permanece em auditoria.

\subsection{Nota de procedência}

Todas as referências citadas neste documento foram conferidas em fonte primária:
depósito do editor, PDF do artigo, amostra oficial da norma ou folha de dados do
fabricante. Foram examinados \num{39} registros. Destes, \num{17} foram confirmados sem
reparo, \num{21} exigiram correção antes de poderem ser citados, e \num{1} foi descartado
por não ter sido possível verificar o achado que se pretendia dele extrair. A
bibliografia da Seção~\ref{sec:bib} tem, portanto, \num{38} itens.

As correções não foram cosméticas, e registrá-las é rigor, não fraqueza. Elas se
distribuem em cinco tipos. Primeiro, listas de autores truncadas por extração automática:
sete registros perdiam de um a quatro autores no fim da lista. Segundo, versão errada do
trabalho: dois artigos estavam datados pelo \emph{preprint} arXiv enquanto se citava o
título da versão de periódico, o que mistura dois documentos distintos. Terceiro, achado
mal atribuído: em três casos o número que se pretendia citar pertence a outro método, a
outro artigo ou a outra edição do mesmo desafio. O caso mais grave é um protocolo de
limiar por distribuição gama que não existe no artigo do DCASE2020 e só aparece na edição
de 2022. Quarto, folhas de dados citadas por revisão e ano errados, o que importa porque
revisões alteram bandas e faixas térmicas. Quinto, e mais consequente para o
Eixo~4: a ISO 10816-3:2009 está \textbf{retirada}, substituída pela ISO 20816-3:2022, e o
levantamento original a listava como vigente; além disso, a URL registrada para a norma
substituta apontava para outra norma inteiramente.

A redação seguiu três regras. Primeira: nenhuma afirmação que dependesse de um número não
confirmado sobreviveu na forma em que estava. Cada uma foi reescrita para o que a fonte
sustenta, ou foi removida. Segunda: valores que vêm de folha de dados do fabricante
ou das notas internas de leitura do projeto, e não da literatura revisada por pares, estão
marcados como tal no corpo do texto e \emph{não} recebem entrada na bibliografia; é o caso
da densidade de ruído do MPU6050 e do orçamento de ruído da Analog Devices discutido no
Eixo~4. Terceira: resultados próprios do Kaelix estão sempre identificados como próprios,
com a indicação de se vêm de simulação, de decimação controlada do MAFAULDA ou de medição
no dispositivo. A distinção importa, porque a camada de aquisição do firmware ainda
não está implementada e, portanto, \emph{nenhum} número deste documento vem de medição no
dispositivo.

Duas normas que o levantamento do projeto cita como pertinentes não foram conferidas em
fonte primária nesta auditoria: a ISO 13373-1 e -2, sobre monitoramento de condição por
vibração, e a ISO 5348, sobre montagem mecânica de acelerômetros. Elas são mencionadas no
texto onde a omissão seria desonesta, sempre com essa ressalva explícita, e não integram a
bibliografia.

\section{Revisão por eixo}

\subsection{Eixo 1 --- Isolation Forest, limiar de decisão e TinyML embarcado}

A literatura recente de detecção não supervisionada de anomalias desloca a pergunta de
engenharia, mas o deslocamento é condicional, e a condição precisa ficar à vista:
\emph{dada banda de medição suficiente}, o que decide o desempenho não é qual detector se
escolhe, e sim como o limiar é calibrado, com que métrica o resultado é lido e quanto do
orçamento de um microcontrolador o modelo consome. Sem banda, nada disso importa. Os
Eixos~3 e~4 mostram que essa é a situação do Kaelix para a falha de rolamento,
onde nenhum corte de escore transforma contraste ausente em detecção. Os grandes
comparativos invocados a seguir são de \emph{benchmark tabular} \emph{offline}, não de
desempenho em campo, e a transferência de um para o outro é hipótese, não resultado.

\subsubsection{Escolha do detector e ausência de consenso na literatura}

Dois comparativos de grande porte discordam, e a discordância é informativa porque eles
medem coisas diferentes. O ADBench é o maior em número de experimentos e em cobertura de
paradigmas: \num{30} algoritmos (\num{14} não supervisionados, \num{7}
semissupervisionados e \num{9} supervisionados) sobre \num{57} conjuntos de dados,
\num{98436} experimentos ao todo, e conclui que \emph{nenhum} método não supervisionado é
estatisticamente superior aos demais~\cite{han2022}. O trabalho de Bouman, Bukhsh e Heskes
é a maior comparação restrita a algoritmos \emph{não supervisionados}: \num{33} algoritmos
sobre \num{52} conjuntos tabulares reais. Nele o Extended Isolation Forest (EIF) supera a
maioria dos concorrentes, com a ressalva de que os dados se separam em dois regimes: em
conjuntos \emph{locais} vence o kNN, em conjuntos \emph{globais} o EIF~\cite{bouman2024}.
Não há contradição entre os dois resultados; há escopos distintos. Mas há consequência
direta para este projeto: a evidência disponível recomenda o EIF, não o Isolation Forest
canônico adotado aqui.

Essa objeção não pode ser apenas nomeada. A razão declarada para permanecer no Isolation
Forest é de implementação embarcada, e é verificável: um nó do iForest canônico armazena
um índice de atributo e um limiar escalar, e a inferência custa uma comparação; um nó do
EIF armazena um vetor normal com tantas componentes quantos forem os atributos, mais um
intercepto, e a inferência custa um produto interno por nó. Com quatro descritores, o
custo por nó cresce de duas para cinco palavras e a inferência passa de comparação para
soma de produtos, em uma árvore percorrida uma vez por janela. É um preço pequeno em
valor absoluto e grande em relação ao orçamento de um ESP32-S3 quando multiplicado pelo
tamanho da floresta. Esse, porém, é um argumento de projeto, e não um resultado de
literatura. Enquanto não for medido no dispositivo, a escolha do iForest canônico deve
constar como limitação de escopo, não como conclusão.

O suporte específico para máquinas rotativas, em contrapartida, é o melhor que o Isolation
Forest tem. Entre \num{11} detectores não supervisionados, ele atinge F1 \num{97,19} e
PR-AUC \num{99,92} em rolamento \emph{run-to-failure}, \num{99,71} e \num{99,99} em caixa
de engrenagens e \num{97,20} e \num{99,74} em faltas mecânicas~\cite{brito2022}. Os dois
primeiros conjuntos são adquiridos a \SI{20}{\kilo\hertz}, banda que o MPU6050 do Kaelix,
com Nyquist de \SI{500}{\hertz}, não possui; o terceiro é descrito de outra forma no
artigo, como arquivos de \num{3200} pontos com resolução espectral de \SI{0,125}{\hertz},
e o registro de auditoria não permite identificar o conjunto de dados nem a taxa de
aquisição correspondente.

Esse terceiro caso merece parágrafo próprio, porque é onde a literatura mais favorável ao
Isolation Forest colide com o resultado próprio deste projeto. Em faltas mecânicas de
bancada, Brito \emph{et al.} reportam F1 \num{97,20}; o Kaelix, com o mesmo detector
sobre o MAFAULDA, mede detecção de defeito de rolamento puro entre \num{15} e
\SI{21}{\percent} a \SI{2}{\percent} de falso alarme. Os dois números não podem coexistir
sem explicação, e três candidatos a explicação estão disponíveis. O primeiro é o limiar:
no artigo, o corte do Caso~3 foi definido a partir da razão de contaminação
\emph{conhecida}~\cite{brito2022}, isto é, com informação que o operador em campo não
tem; um limiar-oráculo desloca o ponto de operação para onde o F1 é máximo, e não para
onde o falso alarme é tolerável. O segundo é o protocolo de partição, que o registro
verificado não documenta. O achado próprio deste projeto é que a partição aleatória de
janelas infla a acurácia em cerca de 36 pontos percentuais em relação à validação por
ensaio. O terceiro só se aplica se o conjunto do Caso~3 for o MAFAULDA, o que não foi
possível confirmar: nesse conjunto, os ensaios de rolamento rodam \emph{com} massa de
desbalanceamento adicionada, e a detecção acompanha a massa e não o defeito, de modo que
reportar sem estratificar por massa infla o resultado. Verificar qual dos três opera, e em
que medida, é trabalho a fazer; enquanto não for feito, o F1 \num{97,20} não pode ser lido
como evidência de que o Isolation Forest resolve rolamento na banda do Kaelix.

Um quarto trabalho é frequentemente invocado como contraponto e não sustenta o que dele
se pede. Hermansa \emph{et al.} comparam quatro detectores (HDBSCAN, LOF, Isolation
Forest e \emph{one-class} SVM) em sensores de britadores de carvão e de pórticos, e
adotam o HDBSCAN como modelo-base do \emph{pipeline} de redução de falsos
alarmes~\cite{hermansa2022}. O registro verificado não contém ranqueamento que coloque o
Isolation Forest atrás dos demais, e essa afirmação foi removida daqui. O que resta é
mais modesto e ainda útil: em uma aplicação industrial real, a escolha do método-base
recaiu sobre outro detector.

\subsubsection{O limiar: o \emph{default} não é calibração}

No artigo fundador, o escore $s(x,n) = 2^{-E(h(x))/c(n)}$ tem $0,5$ como \emph{referência
interpretativa}, na formulação ``if all the instances return $s \approx 0,5$, then the
entire sample does not really have any distinct anomaly'', e a robustez demonstrada está
nos parâmetros do \emph{ensemble} ($t = \num{100}$, $\psi = \num{256}$), não no
corte~\cite{liu2008}. A documentação oficial explicita o mecanismo: com
\texttt{contamination='auto'} o \texttt{offset\_} é fixado em $-0,5$ e
\texttt{decision\_function = score\_samples - offset\_}; qualquer valor diferente redefine
apenas o \emph{offset}, nunca o \texttt{score\_samples}~\cite{scikitlearn2026}. Seguem daí
duas consequências mecânicas: o corte padrão é herdado de uma interpretação, não estimado
nos dados, e toda métrica de ranqueamento, ROC-AUC inclusive, é invariante ao parâmetro.

Os trabalhos que reportam bons resultados não usam o padrão. O limiar é fixado em
\num{0,75} sobre o escore normalizado, sob a hipótese de que o treino contém apenas
comportamento normal~\cite{antonini2023}, ou derivado da razão de contaminação conhecida,
ou da mesma hipótese de treino limpo~\cite{brito2022}. A alternativa principiada é a
Teoria de Valores Extremos: SPOT e DSPOT dispensam hipótese sobre a distribuição
\emph{do corpo} dos dados e fixam o limiar a partir do risco $q$, que controla diretamente
o número de falsos positivos, com DSPOT estendendo o método a fluxos com
deriva~\cite{siffer2017}. Convém não superestimar a facilidade. A teoria não dispensa
hipótese: ela assume que a \emph{cauda} das excedências segue uma Pareto generalizada, o
que exige excedências em número suficiente para estimar seus parâmetros. A implementação
também precisa de mais do que $q$: é preciso escolher o limiar inicial de
\emph{peaks-over-threshold} e a janela de calibração. Além disso, $q$ é uma probabilidade
de excedência sob a cauda ajustada, e não o falso alarme empírico do detector implantado;
os dois coincidem apenas na medida em que o ajuste esteja correto. Nada no material
revisado indica que SPOT ou DSPOT já tenham sido aplicados a escores de Isolation Forest,
nem que o custo caiba na RAM de um ESP32-S3: aplicá-los aqui é proposta deste trabalho,
não resultado herdado da literatura.

Contra tudo isto pode-se invocar o estudo de seleção de modelos \emph{label-free}, que
mostra que nenhuma estratégia interna é praticamente útil: as isoladas não diferem de
sorteio, e as por consenso não superam o iForest com hiperparâmetros padrão e ainda custam
mais~\cite{ma2023}. Poder-se-ia concluir daí que o valor padrão equivale a qualquer outra escolha.
Não se pode: o registro verificado descreve \num{297} modelos candidatos formados por
\num{8} detectores e várias configurações de hiperparâmetros, sem enumerá-las, e não
reporta \texttt{contamination} entre os hiperparâmetros varridos. Como o parâmetro não
altera \texttt{score\_samples}, ele não afeta o ranqueamento sobre o qual aquele estudo se
pronuncia. A conclusão de Ma \emph{et al.}, de que não existe método confiável sem rótulos
para escolher detector e hiperparâmetro, permanece válida e não se transfere ao limiar.

\subsubsection{Custo em microcontrolador, métricas e falso alarme}

Em implementação embarcada do Isolation Forest, o trabalho mais próximo do Kaelix treina e
executa o detector no próprio ESP32, em MicroPython, com IMU amostrada a
\SI{1125}{\hertz}~\cite{antonini2023}. Os números verificados são \SI{84}{\kilo\byte}
(\num{84815}~B) de memória na inferência com \num{50} árvores, um pico de
\SI{1570}{\kilo\byte} no treino, cerca de \SI{1}{\mega\byte} de \emph{heap} para carregar
o modelo de \num{50} árvores (arquivo de \SI{284}{\kilo\byte}), latência de
\SI{20,33}{\milli\second} ($\pm\,\SI{7,17}{\milli\second}$) em instâncias normais e
\SI{16,27}{\milli\second} ($\pm\,\SI{7,04}{\milli\second}$) nas anômalas, e treino
\emph{on-device} de \SI{1,2}{\second} com \num{10} árvores a \SI{6,4}{\second} com
\num{50}. Duas ressalvas acompanham esses valores e valem para qualquer uso deles como
orçamento. A primeira é de atribuição: o registro de auditoria associa o pico de
\SI{1570}{\kilo\byte} ao treino com \num{30} árvores, enquanto as notas de leitura do
projeto registram uma faixa de \SIrange{536}{1576}{\kilo\byte} para \numrange{10}{50}
árvores; a divergência precisa ser resolvida contra o artigo antes de o número ser
transportado. A segunda é de natureza: trata-se de implementação em MicroPython
interpretado, e o Kaelix é C++ nativo. Comparar pegadas de memória e latências entre as
duas sem dizê-lo é comparar coisas diferentes. O orçamento correspondente do Kaelix ainda
não foi medido, porque a camada de aquisição do firmware permanece não implementada, e
por isso este documento não apresenta números próprios de RAM e de latência.

O único trabalho que ataca o Isolation Forest especificamente para TinyML argumenta que
``the standard algorithm might be improved in terms of memory requirements, latency and
performances'' e usa supervisão fraca para podar árvores, operando com \numrange{20}{40}
em vez das \num{100} padrão~\cite{barbariol2022}. É a via de redução de custo mais direta
disponível, e ela exige alguns rótulos.

Quanto à métrica, o benchmark de referência para o mesmo \emph{regime de treino} do Kaelix
--- apenas dados normais no treino, anomalias desconhecidas --- adota AUC \emph{e} pAUC
sobre a faixa de baixa FPR $[0, p]$ com $p = \num{0,1}$, porque a AUC global não
caracteriza o ponto de operação~\cite{koizumi2020}. A identidade com o caso do Kaelix
termina aí: aquele é um desafio de detecção de som anômalo de máquina, com modalidade,
banda e classe de modelo distintas das deste projeto. A justificativa do pAUC não depende
dessa aproximação, e Brito \emph{et al.} reforçam a mesma disciplina ao reportar PR-AUC
além do F1~\cite{brito2022}. Isso ampara a decisão de reportar pAUC a FPR $\leq
\SI{10}{\percent}$, e um resultado próprio mostra por que ela importa: em \emph{sinal
sintetizado}, comparando a cadeia de referência a \SI{20}{\kilo\hertz} com a do MPU6050,
a pAUC cai \SI{16,2}{\percent} onde a AUC global cai apenas \SI{3,7}{\percent}. O número
não é de dado medido, e não é comparável aos resultados sobre o MAFAULDA discutidos
adiante.

Sobre falso alarme, a evidência mais direta e mais desfavorável ao desenho do Kaelix é a
de Chevtchenko \emph{et al.}: em nó IoT com ESP32 e IMU ISM330DLCTR amostrada a
\SI{1}{\kilo\hertz}, sobre \num{3750} amostras normais e \num{2250} anômalas de três
motores de indução de 1~hp a \num{1000}~rpm, o Isolation Forest atinge sensibilidade de
\SI{88,9}{\percent} em validação e \SI{86,3}{\percent} em teste, mas especificidade de
apenas \SI{61,6}{\percent} e \SI{67,4}{\percent}, isto é, de \SI{32,6}{\percent} a
\SI{38,4}{\percent} da operação normal marcada como anomalia~\cite{chevtchenko2023}.
Existe, portanto, FPR publicado para Isolation Forest embarcado, e ele é alto. A
evidência de campo sobre redução de falso alarme com detector não supervisionado genérico
é complementar e mais otimista: a métrica $N$ cai de \numrange{0,457}{0,939} para
\numrange{0,026}{0,115} com adaptação do modelo-base somada a um modelo de correção, uma
redução média de \SI{90,25}{\percent}~\cite{hermansa2022}. Esses números, porém, são do
HDBSCAN, o modelo-base escolhido, e não do Isolation Forest.

A Tabela~\ref{tab:fpr} reúne o que se sabe. Os \SI{22}{\percent} atribuídos a
Kolok \emph{et al.} não são reportados no artigo: são derivados aqui a partir dos
valores publicados de acurácia (\num{0,729}), precisão (\num{0,755}) e revocação
(\num{0,679}), que juntos determinam prevalência de \num{0,50} e especificidade de
\num{0,78}. A derivação é aritmética e só vale se os três valores forem internamente
consistentes. O artigo é internamente inconsistente em outro ponto, ao mencionar
\SI{93,4}{\percent} de acurácia diagnóstica na discussão sem que o número apareça nos
resultados. Ainda assim, o quadro é o seguinte: três implementações independentes com
Isolation Forest em nó embarcado ficam entre \num{22} e \SI{38}{\percent} de falso alarme
sobre operação normal, e o valor próprio de cerca de \SI{42}{\percent} com
\texttt{contamination='auto'} é da mesma ordem de grandeza. É preciso dizer o que essa
compatibilidade \emph{não} é: Chevtchenko \emph{et al.} não declaram, no registro
verificado, qual política de limiar usaram, de modo que a proximidade numérica entre
\SI{38}{\percent} e \SI{42}{\percent} é compatibilidade de ordem de grandeza entre
conjuntos, máquinas e limiares diferentes, e não um correlato empírico controlado. E há
uma assimetria mais forte que essa: os três valores da literatura são medidos em máquinas
reais, ao passo que os \SI{42}{\percent} próprios vêm de sinal sintetizado, com detecção
de \SI{100}{\percent} e AUC unitária --- condições que nenhum conjunto real produz. A
compatibilidade sustenta que \texttt{contamination='auto'} falha por construção, e não
que o Kaelix reproduza, em máquina, o falso alarme dessas implementações.

\begin{table}[!htb]
\centering
\caption{\textbf{Falso alarme de Isolation Forest em nó embarcado.} Fração da operação
normal marcada como anomalia. O valor de Kolok \emph{et al.} é derivado neste documento
a partir de acurácia, precisão e revocação publicadas; os demais são reportados pelas
fontes. Os dois últimos são resultado próprio, e têm procedências distintas: o valor com
\texttt{contamination='auto'} vem de \emph{sinal sintetizado}, e o valor calibrado vem do
MAFAULDA com validação \emph{GroupKFold} por ensaio. Apenas as três primeiras linhas e a
última correspondem a medição sobre máquina.}
\label{tab:fpr}
\footnotesize
\setlength{\tabcolsep}{4pt}
\begin{tabular}{@{}p{3.1cm}p{4.0cm}p{4.1cm}p{3.1cm}@{}}
\toprule
\textbf{Trabalho} & \textbf{Plataforma e sensor} & \textbf{Política de limiar} & \textbf{Falso alarme} \\
\midrule
Kolok \emph{et al.}~\cite{kolok2025}
  & ESP32-C6 + MPU6050
  & \texttt{contamination} \texttt{='auto'}
  & $\approx\SI{22}{\percent}$ (derivado) \\
Chevtchenko \emph{et al.}~\cite{chevtchenko2023}
  & ESP32 + ISM330DLCTR a \SI{1}{\kilo\hertz}
  & não declarada no registro verificado
  & \SIrange{32,6}{38,4}{\percent} \\
Antonini \emph{et al.}~\cite{antonini2023}
  & ESP32-WROVER-IE + IMU a \SI{1125}{\hertz}
  & escore normalizado fixo em \num{0,75}
  & não reportado \\
Kaelix (próprio, \emph{sinal sintetizado})
  & ESP32-S3 + MPU6050
  & \texttt{contamination} \texttt{='auto'}
  & $\approx\SI{42}{\percent}$ \\
Kaelix (próprio)
  & ESP32-S3 + MPU6050
  & calibrado para FPR alvo
  & \SI{2}{\percent} (ponto adotado) \\
\bottomrule
\end{tabular}
\end{table}

Restam duas objeções de peso equivalente ao das convergências. A primeira: com apenas
\SI{1}{\percent} de anomalias rotuladas a maioria dos métodos semissupervisionados supera
o melhor não supervisionado~\cite{han2022}, o que torna a aposta puramente não
supervisionada um custo deliberado num conjunto rotulado como o MAFAULDA. A resposta
declarada é de cenário de implantação, e não de desempenho: o dispositivo é destinado a
máquinas sem histórico rotulado, os rótulos do MAFAULDA são usados aqui para
\emph{validar} e não para treinar, e um detector que só funciona com rótulos não se
implanta em uma máquina nova. Onde esse critério não se aplicar, como em uma planta que já
tenha histórico anotado, a escolha não se defende e deve ser tratada como limitação de
escopo. A segunda objeção corre na direção contrária: trocar o Isolation Forest por um
autoencoder não resolve, pois anomalias distantes dos dados normais podem ser
perfeitamente reconstruídas~\cite{bouman2025} --- com a ressalva de que essa fonte é
\emph{preprint} sem revisão por pares.

Duas lacunas ficam abertas neste eixo, e uma terceira foi eliminada pela auditoria. A que
caiu era a afirmação de que não existiria FPR publicado para Isolation Forest embarcado:
existe, e está na Tabela~\ref{tab:fpr}. O que sobrevive é mais estreito e mais defensável.
Primeiro, ninguém isolou o efeito do parâmetro: não se localizou trabalho que, sobre a
mesma base e o mesmo modelo, varie \texttt{contamination} e reporte a fração de operação
normal marcada, em vibração industrial de banda limitada; Kolok \emph{et al.} usam
\texttt{'auto'} sem variá-lo, Chevtchenko \emph{et al.} não declaram a política e
Antonini \emph{et al.} fixam \num{0,75} e declaram que ``an evaluation of the detection
accuracy is beyond the scope of this work''~\cite{antonini2023}. Segundo, não há falso
alarme reportado para Isolation Forest \emph{treinado no próprio dispositivo}: o único
trabalho que treina \emph{on-device} é justamente o que declara a acurácia fora de escopo.
E ninguém cruza as três restrições simultâneas deste projeto: acelerômetro MEMS limitado
a \SI{500}{\hertz}, detector não supervisionado embarcado e limiar calibrado para um FPR
alvo.

\subsection{Eixo 3 --- Banda de acelerômetros MEMS e detectabilidade de falha de rolamento}

A tensão que organiza esta literatura é uma só: as frequências que se deseja medir num
rolamento são baixas, mas o mecanismo físico que as torna mensuráveis é alto. Um defeito
localizado produz impactos que excitam ressonâncias estruturais, e a taxa de repetição do
defeito aparece como \emph{modulação} dessa portadora, não como energia própria no
espectro direto; daí o método de referência ser a análise de envelope, que filtra o sinal
numa banda alta antes de demodular~\cite{randall2011}. Landi \emph{et al.} tornam o
contraste explícito: no rolamento 22209E1K a \num{3000}~rpm, BPFO, BPFI, BSF e FTF valem
\SI{361}{\hertz}, \SI{488}{\hertz}, \SI{161}{\hertz} e \SI{21}{\hertz}, todas abaixo de
\SI{500}{\hertz}; ainda assim a emulação de falha adota ressonância dos componentes em
\SI{10}{\kilo\hertz} e amostragem de \SI{160}{\kilo\hertz}, e o sensor julgado adequado é
um cubo de três ADXL1005Z cobrindo \SIrange{10}{25000}{\hertz}, com banda de exatidão de
\SI{10}{\percent} até \SI{12}{\kilo\hertz}~\cite{landi2023}. O alvo cabe na banda do
Kaelix; o mecanismo de detecção, não.

A mesma geometria explica por que descritores estatísticos calculados na banda base
falham. A curtose não é um número global útil, e sim uma \emph{função} da frequência, que
indica em qual banda os transientes do defeito são impulsivos~\cite{antoni2006}; o
kurtograma rápido formaliza a busca da díade ótima --- frequência central e largura de
banda --- que maximiza essa impulsividade, para escolher o passa-banda da
demodulação~\cite{antoni2007}. Um dispositivo cujo Nyquist é \SI{500}{\hertz} não tem
díade a escolher: mede a impulsividade de uma banda onde o defeito quase não deposita
energia.

\subsubsection{De onde vêm os números próprios deste eixo}

Antes de confrontá-los com a literatura, é preciso dizer o que os resultados de banda do
Kaelix são e o que não são. Eles têm duas procedências distintas, e tratá-las como uma só
seria cometer aqui o defeito que este eixo cobra da literatura. Os resultados de
detecção vêm de decimação controlada, em \emph{software}, dos sinais do MAFAULDA
adquiridos a \SI{50}{\kilo\hertz}. Já o par \num{413}$\times$/\num{20}$\times$ de contraste
em BPFO não vem do MAFAULDA: vem de sinal \emph{sintetizado} a partir de modelo
analítico --- trem de impulsos amortecidos modulando uma ressonância estrutural ---
comparado entre a cadeia de referência a \SI{20}{\kilo\hertz} e a cadeia do MPU6050 a
\SI{1}{\kilo\hertz}. \emph{Nenhum} dos dois vem de medição no dispositivo. A distinção não é formalidade. O firmware ainda devolve
\texttt{Status::NotImplemented} na inicialização do acelerômetro, e o filtro passa-baixas
interno do MPU6050 permanece como pendência declarada no código: sem antialiasing
configurado abaixo de $f_s/2$, todo conteúdo acima de \SI{500}{\hertz} dobra para dentro
da banda de análise e fica indistinguível de sinal legítimo. Este eixo cobra da literatura
precisamente ``antialiasing controlado'', e seria desonesto não registrar que o próprio
dispositivo ainda não o tem.

A consequência é um artefato que merece registro por ser, ele próprio, um achado. Sem
DLPF, uma ressonância de alguns quilohertz não desaparece: ela reaparece rebatida em
alguma frequência dentro da banda, e a curtose e o fator de crista medidos sobre esse
sinal podem \emph{parecer} separar as classes por uma razão que nada tem a ver com a
física do defeito. Uma medição ingênua com MPU6050 pode, portanto, funcionar pelos motivos
errados e passar em validação enquanto o mecanismo real é aliasing. Não se localizou,
na literatura revisada por pares de diagnóstico de rolamento, descrição experimental desse
artefato; ele aparece em notas de aplicação de fabricante, registradas no levantamento do
projeto e não conferidas nesta auditoria.

Com essa marcação feita, o resultado próprio é: sobre o mesmo sinal decimado, a taxa de
repetição do defeito permanece dentro da banda, mas o contraste do indicador cai de
\num{413} para \num{20} vezes, e curtose e fator de crista deixam de separar as classes.
Duas qualificações se impõem. A primeira é que ``a taxa de repetição está dentro da
banda'' não é o mesmo que ``a taxa de repetição aparece no espectro'': localizou-se na
busca um trabalho (Hu, Xiang, Xu e Lin, 2019, \emph{Chinese Journal of Mechanical
Engineering}) cujo resumo sustenta que BPFO e BPFI podem \emph{desaparecer} do espectro
por cancelamento de fase entre múltiplas respostas ressonantes, mesmo com banda adequada;
como o texto integral não pôde ser aberto, a referência não integra esta bibliografia e a
objeção fica registrada como pendência: verificar, no MAFAULDA decimado, em que fração dos
ensaios o pico de BPFO é de fato observável. A segunda é que o par
\num{413}$\times$/\num{20}$\times$ não mede o que se pretendia: \num{413} é envelope na
banda de ressonância a \SI{20}{\kilo\hertz} e \num{20} é espectro direto a
\SI{1}{\kilo\hertz}: mudam duas variáveis --- a banda e o método de demodulação --- e o
efeito estava sendo atribuído a uma. Não é ``a mesma métrica'', e dois pontos não são
``função da banda''.

\subsubsection{Quão estreita pode ser a banda de um MEMS}

Os estudos de adequação convergem para o limite superior de banda como variável determinante,
mas nenhum chega perto de \SI{500}{\hertz}. Albarbar \emph{et al.} compararam três MEMS
capacitivos (\SIrange{1}{6000}{\hertz}, \SIrange{1}{10000}{\hertz} e \SI{1500}{\hertz}) a
um piezo ICP de \SIrange{1}{2000}{\hertz} e \SI{100}{\milli\volt\per\gram}: o de banda
mais estreita foi declarado inadequado para monitoramento de condição~\cite{albarbar2008}.
A razão não é a banda: ele foi reprovado por sensibilidade instável
(\SIrange{63}{111}{\milli\volt\per\gram}), ruído e desvio de fase substanciais. Mesmo os
dois de banda larga produziram, em máquina CNC, ``\textit{lots of noise including extra
un-interpretable peaks}''~\cite{albarbar2008}. O dado de sensibilidade é usado aqui apenas
de forma qualitativa: o registro verificado documenta a faixa medida e o piezo de
referência, mas a rotulagem dos sensores no artigo é internamente inconsistente, e a
sensibilidade nominal com que a faixa medida deveria ser comparada não consta do registro
(ela aparece nas notas do projeto como valor de folha de dados). O Eixo~4 retoma o ponto.

Fidali \emph{et al.} avaliaram dois sensores MEMS digitais contra referência
piezoelétrica em onze rolamentos 6303 de esferas com gaiola polimérica, seis novos e
cinco alterados: pista externa, pista interna, elemento rolante, gaiola progressivamente
danificada e, como quinta condição, ausência de lubrificante, que não é defeito mecânico e
sim condição de lubrificação~\cite{fidali2024}. Duas grandezas estavam sendo confundidas
neste documento: as faixas de catálogo dos sensores são
\SIrange{2}{2500}{\hertz} e \SIrange{10}{4000}{\hertz}, enquanto as bandas em que os
descritores foram efetivamente calculados são \SIrange{2}{3500}{\hertz} e
\SIrange{1000}{4000}{\hertz}. Uma é especificação de sensor, a outra é banda de análise, e
o Eixo~4 usa a mesma leitura. Os resultados: a curtose foi bem menos sensível nos MEMS que
no piezo, e para \emph{todos} os descritores a menor sensibilidade foi ao dano de gaiola.
Os autores relatam que a sensibilidade diagnóstica dos parâmetros depende do limite
superior de banda dos sensores, mas registram no mesmo texto que ``\textit{at this stage
of the research, it is not possible to state unequivocally that the use of a piezoelectric
transducer increases the diagnostic sensitivity to a decisive degree compared to MEMS
capacitive accelerometers}''~\cite{fidali2024}. Esta é a leitura adotada nos dois eixos:
tendência observada, não veredito. E a ressalva dos autores diz respeito a MEMS de
\SIrange{2,5}{4}{\kilo\hertz}, não a um sensor de \SI{500}{\hertz}.

O que se pode afirmar sem esticar as fontes é, portanto, um enunciado positivo e mais
fraco do que o que estava escrito: \emph{nenhum estudo de adequação revisado testou sensor
com limite superior abaixo de \SI{1500}{\hertz}}. Comparar isso com o Kaelix exige
comparar banda com banda, e não banda com Nyquist: o MPU6050 a \SI{1}{\kilo\hertz} de taxa
de saída tem banda útil de, no máximo, \SI{500}{\hertz}, e isso somente se o DLPF for
configurado abaixo de $f_s/2$, o que ainda não foi feito. Contra os \SI{1500}{\hertz}
do sensor mais estreito já submetido a um estudo de adequação, é um terço da banda. E o
argumento por ausência não vale: ninguém ter reprovado algo mais estreito não é o mesmo
que alguém ter aprovado.

A evidência arquitetural reforça a leitura. Jakobsen, projetando um nó IoT
com LoRa e bateria próximo do Kaelix, separou os papéis entre dois acelerômetros: o
ADXL1002 (\num{25}~\dens, resposta plana DC--\SI{11}{\kilo\hertz}, ressonância em
\SI{21}{\kilo\hertz}, amostrado a \SI{40}{\kilo\hertz}) para falha de rolamento,
confirmada por $F_o$ e harmônicos no espectro de envelope em \SIrange{0}{700}{\hertz}, e o
ADXL362 (\num{550}~\dens, amostrado a \SI{400}{\hertz}) apenas para desbalanceamento e
rotação abaixo de \SI{200}{\hertz}~\cite{jakobsen2024}. Sharma \emph{et al.} repetem a
escolha em bancada de rolamentos e engrenagens com ADXL1002Z e \textit{Random
Forest}~\cite{sharma2024}. O MPU6050 do Kaelix tem densidade de ruído da ordem de
\num{400}~\dens{}, valor de folha de dados do fabricante (InvenSense), não conferido
experimentalmente e, por isso, sem entrada nesta bibliografia. Se ele estiver correto, o
sensor está mais próximo daquele que Jakobsen restringiu a
\SI{200}{\hertz} do que daquele que ele reservou ao rolamento. A ressalva é a mesma que o
Eixo~4 faz à ausência de calibração rastreável: um número de fabricante não é uma medição.

\subsubsection{As objeções}

A contra-evidência é real e merece o mesmo peso, mas cada objeção precisa ser respondida
pelo que ela é. El Bouharrouti \emph{et al.} reamostraram sinais do CWRU de
\SI{48}{\kilo\hertz} até \SI{1}{\kilo\hertz}, em passos de \SI{6}{\kilo\hertz} mais o
ponto de \SI{1}{\kilo\hertz}, sobre dez estados de saúde --- um saudável e nove
defeituosos, com esfera, pista interna e pista externa em três severidades cada --- e
concluem que ``\textit{depending on the machine learning algorithm used, better training
accuracies are not systematically achieved when training with vibration signals sampled at
a relatively high frequency}''~\cite{elbouharrouti2024}. É classificação supervisionada
(SVM, MNLR, MLP, LSTM e XGBoost) com \textit{downsampling} fracionário correto, e não
um IMU com anti-\textit{aliasing} de prateleira. O registro verificado não documenta o
protocolo de partição, de modo que não é possível afirmar que o vazamento entre janelas
tenha sido controlado; verificá-lo é pendência, não acusação.

Maruthi e Hegde detectaram múltiplas falhas de rolamento em motor de indução por
FFT do espectro \emph{direto} com acelerômetro MEMS, sob condições forçadas de operação
--- sem carga, monofasamento e tensão desbalanceada ---, com bandas laterais em torno da
fundamental indicando modulação do entreferro e frequências experimentais próximas das
analíticas~\cite{maruthi2016}. Aqui não há classificador treinado, e portanto a crítica de
vazamento por ensaio é categoria inaplicável. São falhas semeadas, de severidade não incipiente, sob condições elétricas que amplificam a
assinatura.

O caso mais próximo do Kaelix \emph{em hardware e em modo de falha} é o mais sóbrio. Kolok
\emph{et al.} construíram protótipo com MPU6050 e microfone MEMS INMP441 sobre MCU da
mesma classe, com Isolation Forest não supervisionado treinado só em dados saudáveis, e
obtiveram \SI{72,9}{\percent} ($\pm$~\SI{2,4}{\percent}) de acurácia em validação cruzada
de cinco dobras, com revocação de \num{0,679} --- para uma única falha, obstrução
deliberada de roda, isto é, desbalanceamento; o desgaste de rolamento foi explicitamente
adiado para trabalho futuro, e o \textit{streaming} efetivo foi de uma linha a cada
\SI{0,5}{\second}, cerca de \SI{2}{\hertz}, contra taxa interna de
\SI{1}{\kilo\hertz}~\cite{kolok2025}. Um dado adicional, registrado nas notas de leitura
do projeto, é que o filtro passa-baixas interno do MPU6050 foi configurado em
\SI{5}{\hertz}. Ele não consta do registro de auditoria e precisa ser reconfirmado no
artigo antes de sustentar comparação. Se assim for, a banda efetiva daquele trabalho está
duas ordens de grandeza abaixo da do Kaelix, e ele não é evidência sobre o que um MPU6050
com Nyquist de \SI{500}{\hertz} consegue: é evidência sobre desbalanceamento em banda
quase contínua. O padrão qualitativo de que o MPU6050 resolve desbalanceamento e cala
sobre rolamento é o que os resultados do Kaelix em dados reais reproduzem, mas o número
\SI{72,9}{\percent} não pode ser lido como teto da classe.

\subsubsection{As alternativas que não foram tomadas}

Este eixo trata \SI{500}{\hertz} como restrição, mas ela é escolha de
componente, não limite de classe. Existem IMUs de faixa de preço equivalente com taxa de
saída muito superior à do MPU6050; esta revisão não localizou, no conjunto auditado,
nenhum estudo que as caracterize em falha de rolamento, de modo que a comparação fica
registrada como decisão de projeto a documentar por folha de dados, e não como
resultado de literatura. A permanência no MPU6050 precisa então ser justificada por
critério declarado: disponibilidade do componente, custo e escopo do trabalho. Deixá-la
implícita produz a impressão de restrição física onde há escolha.

A segunda alternativa é arquitetural e tem apoio direto nas fontes deste documento: três
dos nós revisados recuperam a banda alta por canal acústico em vez de acelerômetro. O nó
de Jakobsen combina acelerômetro MEMS e microfone já no título do
trabalho~\cite{jakobsen2024}; o protótipo de Kolok \emph{et al.} inclui um microfone
MEMS INMP441 ao lado do MPU6050~\cite{kolok2025}; e o nó de Chevtchenko \emph{et al.}
traz dois microfones SPH0645LM4H-B amostrados a \SI{20}{\kilo\hertz}, ao lado da IMU de
\SI{1}{\kilo\hertz}~\cite{chevtchenko2023}. Um microfone MEMS custa uma fração do
orçamento do dispositivo e alcança a década de frequência que o acelerômetro descarta.
Nenhum dos três, porém, publica resultado de rolamento obtido pelo canal acústico, de modo
que a via está aberta e não está demonstrada.

A contribuição possível deste eixo é, com essas correções, mais estreita e mais honesta do
que a que se pretendia. Não é ``quantificar a degradação do contraste diagnóstico em
função da banda'', porque o par medido muda duas variáveis ao mesmo tempo. É executar a
varredura com método fixo --- envelope em todos os pontos, ou espectro direto em todos os
pontos --- sobre a mesma gravação decimada em vários $f_s$, e reportar a decomposição:
quanto da queda de contraste é banda e quanto é perda do envelope. Essa varredura a
literatura só discute qualitativamente, e ela é executável com os dados já disponíveis.

\subsection{Eixo 4 --- Normalização ISO e integração aceleração--velocidade}

A avaliação de severidade não é convenção arbitrária de unidades, mas escolha física
justificada no próprio texto normativo. O Anexo~D da ISO 20816-3:2022 afirma que vibrações
com o mesmo valor eficaz de velocidade na faixa de \SIrange{10}{1000}{\hertz} podem, em
geral, ser consideradas de igual severidade, e explicita que, se deslocamento ou
aceleração fossem usados na avaliação, os critérios variariam com a
frequência~\cite{iso208163_2022}. Daí a forma das tabelas de zona: os limites existem em
\si{\milli\meter\per\second} eficazes, no Grupo~2 em fundação rígida,
\SI{1,4}{\milli\meter\per\second} em A/B, \SI{2,8}{\milli\meter\per\second} em B/C e
\SI{4,5}{\milli\meter\per\second} em C/D. Os deslocamentos tabelados (\num{22}, \num{45} e
\SI{71}{\micro\meter}) são \emph{calculados} a partir da velocidade supondo frequência
dominante de \SI{10}{\hertz}, não medidos~\cite{iso208163_2022}.

Antes de qualquer discussão de banda, é preciso declarar o escopo, porque ele decide a
quem a norma se aplica. A ISO 20816-3:2022 vale para conjuntos de máquina com potência
acima de \SI{15}{\kilo\watt} e rotação entre \num{120} e \num{30000}~r/min, e a fronteira
de \SI{1,4}{\milli\meter\per\second} pertence ao Grupo~2, isto é, máquinas de
\SIrange{15}{300}{\kilo\watt} com altura de eixo entre \num{160} e
\SI{315}{\milli\meter}~\cite{iso208163_2022}. A bancada do MAFAULDA, sobre a qual todos os
resultados próprios em dados reais foram obtidos, está muito abaixo desse escopo. Isso não
invalida o uso do conjunto: ele é usado aqui como \emph{sinal}, para estudar a cadeia de
medição. Invalida, sim, tratá-lo como caso de aplicação da norma, e o enquadramento das
máquinas-alvo do Kaelix precisa ser declarado no projeto antes que qualquer rótulo de zona
seja emitido.

A ISO 10816-3:2009, tratada como vigente em parte da literatura aplicada e em parte do
levantamento que originou este documento, foi \textbf{retirada} e substituída pela
ISO 20816-3:2022, que cancela também as emendas de
2017~\cite{iso108163_2009,iso208163_2022}. Os valores de zona são idênticos nas duas, de
modo que a substituição afeta a rastreabilidade da citação, e não os números. A norma
antiga é mantida na bibliografia porque é a que a prática industrial e boa parte da
literatura ainda adotam, e porque é nela que está a delegação explícita à ISO 2954.

\subsubsection{A banda normativa como requisito de instrumentação}

Ambas as edições exigem resposta plana em faixa de \emph{ao menos}
\SIrange{10}{1000}{\hertz} para a medição do valor eficaz de banda larga, com limite
inferior de até \SI{2}{\hertz} em máquinas próximas ou abaixo de
\num{600}~r/min~\cite{iso108163_2009,iso208163_2022}. A edição de 2009 delega esse
``plana'' à ISO 2954, que o torna verificável: tolerância de banda passante de
\SI{+0,83}{\decibel} a \SI{-0,92}{\decibel} --- isto é, \SI{10}{\percent} --- de
\SI{12,6}{\hertz} a \SI{1000}{\hertz}, sensibilidade \emph{nominal} a \SI{10}{\hertz} já
de \SI{-3,01}{\decibel} (\num{0,707}) com tolerância de $\pm$\,\SI{2}{\decibel}, e
limitação de banda por par de Butterworth de terceira ordem com $f_1=\SI{10}{\hertz}$
($Q_1=Q_3=1/2$) e $f_2=\SI{1000}{\hertz}$ ($Q_2=Q_4=1/\sqrt{2}$)~\cite{iso2954_2012}. A
norma impõe ainda disciplina de relato relevante a dispositivos truncados: as
frequências de corte de \SI{3}{\decibel} devem ser declaradas junto ao valor
medido~\cite{iso2954_2012}. Para um MPU6050 a \SI{1}{\kilo\hertz}, com Nyquist em
\SI{500}{\hertz}, isso não é diferença de grau: o corte superior fica uma oitava abaixo do
exigido. A leitura honesta é $v_{\mathrm{rms}}$ declarado em
\SIrange{10}{500}{\hertz}, jamais ``conforme ISO 20816-3''.

Este documento omitia um requisito normativo mais severo, e convém desfazer a omissão
mesmo sem fonte primária para citar. O levantamento do projeto registra que a ISO 13373-1
e -2, que tratam de monitoramento de condição por vibração e cobrem especificamente
rolamentos e engrenagens, recomendam \emph{aceleração} como grandeza para esse tipo de
falha e indicam limite superior da ordem de \SI{10}{\kilo\hertz}, podendo ser aumentado.
Essas normas não foram conferidas em fonte primária nesta auditoria e por isso não constam
da bibliografia; mas, se a leitura estiver correta, a consequência é dura e precisa ser
dita: para a falha-alvo do Kaelix o requisito normativo de instrumentação é violado por
uma ordem de grandeza, e não por uma oitava. Citar a norma que o projeto quase atende e
calar sobre a norma que o reprova seria seleção de evidência.

\subsubsection{Integração: obrigação normativa e problema numérico}

A cláusula 4.3 da ISO 20816-3:2022 é explícita: acelerômetros montados em partes
estacionárias, prática comum, \emph{devem} ter a saída integrada para fornecer sinal
de velocidade, com ressalva quanto à dupla integração e à filtragem
passa-alta~\cite{iso208163_2022}. A norma obriga a integrar, mas não prescreve o método.
Brandt e Brincker ocupam esse vazio: a regra do trapézio deve ser evitada por
produzir resultados enviesados, particularmente em frequências mais altas, enquanto o
percurso DFT $\rightarrow$ divisão por $j\omega$ $\rightarrow$ IDFT é o melhor método para
registros acima de 16K amostras, com filtro IIR otimizado de exatidão comparável, e WOLA
só é aceitável em regime permanente~\cite{brandt2014}. São dois mecanismos distintos: o
viés descrito ali é erro de amplitude dependente da frequência, e não a deriva por
\emph{offset} DC residual que domina a integração de MEMS de baixo custo.

Os erros de \SIrange{73}{716}{\percent} medidos no Kaelix com integração no tempo,
contra \SI{0,0}{\percent} na frequência com máscara de banda, combinam os dois mecanismos.
A forma como esse par foi apresentado precisa ser corrigida em dois pontos. O
primeiro: \SI{0,0}{\percent} não é medição, é identidade. Quando o sinal de teste é
sintetizado como soma de senoides contidas na máscara e a integração é feita na mesma base
DFT, o resultado é exato por construção, e nada diz sobre desempenho com sinal real, onde
há transiente, vazamento espectral e janelamento. O segundo: a faixa de \numrange{73}{716}
foi apresentada sem as condições que a geram --- frequência dominante, \emph{offset} DC,
comprimento da janela, presença de rampa. A omissão impede reprodução e sugere escolha do
pior caso. O que o relatório de simulação de fato fixa é um ponto verificável: com viés de
\SI{0,02}{\meter\per\second\squared}, valor contido na especificação do sensor, a
integração na frequência com máscara devolve o valor de referência, enquanto a integração
no tempo deriva. Reportar isso corretamente exige tabelar os erros do domínio do tempo por
condição, e repetir a integração na frequência sobre sinal real do MAFAULDA, com
janelamento e vazamento presentes.

A lacuna estava enunciada assim: ``não se localizou trabalho que quantifique o erro de
integração para $v_{\mathrm{rms}}$ na banda normativa''. Nessa forma, era construída por
excesso de especificação. Brandt e Brincker quantificam o erro dos métodos de integração;
o que não está publicado é a \emph{aplicação} desses métodos ao indicador da norma, o que
é cálculo, e não lacuna de conhecimento. Há, em compensação, um número decisivo que este
projeto pode computar e ainda não computou: quanto da energia de $v_{\mathrm{rms}}$ de um
motor está entre \num{500} e \SI{1000}{\hertz}, isto é, quanto um dispositivo truncado
\emph{subestima} a severidade. Com o MAFAULDA a \SI{50}{\kilo\hertz} já disponível,
comparar $v_{\mathrm{rms}}(\SIrange{10}{1000}{\hertz})$ com
$v_{\mathrm{rms}}(\SIrange{10}{500}{\hertz})$ é imediato, e transforma em resultado a
afirmação qualitativa ``o truncamento não é diferença de grau''.

\subsubsection{O que a literatura de sensores objeta}

O trabalho mais próximo \emph{na aplicação da banda normativa a MEMS de banda limitada} é
o de Fidali \emph{et al.}, que compara sensores digitais de catálogo
\SIrange{2}{2500}{\hertz} e \SIrange{10}{4000}{\hertz} --- com descritores calculados,
respectivamente, em \SIrange{2}{3500}{\hertz} e \SIrange{1000}{4000}{\hertz} --- contra
referência piezoelétrica de \SIrange{0,5}{10000}{\hertz}~\cite{fidali2024}. Os autores
relatam que a sensibilidade diagnóstica depende do limite superior de banda dos sensores e
que os descritores mais sensíveis à fadiga de rolamento são pico e pico a pico de
\emph{aceleração}, não a velocidade eficaz, o que tensiona o critério normativo, mas
registram, no mesmo texto, que naquele estágio não era possível afirmar inequivocamente
que o piezoelétrico aumentasse de forma decisiva a sensibilidade diagnóstica frente aos
MEMS capacitivos~\cite{fidali2024}. É a mesma leitura adotada no Eixo~3: tendência
observada, não veredito. E a ressalva se refere a MEMS de \SIrange{2,5}{4}{\kilo\hertz},
não a um sensor de \SI{500}{\hertz}.

A própria norma reconhece a cisão. A NOTA da cláusula 3.2 da ISO 10816-3:2009 já advertia
que, para fins diagnósticos, um limite superior acima de \SI{1000}{\hertz} pode ser
necessário~\cite{iso108163_2009}. Severidade e diagnóstico são perguntas distintas, e
nenhuma delas é respondida por um sensor limitado a \SI{500}{\hertz}.

Há objeção anterior à banda, e mais desconfortável: a escala. Albarbar \emph{et al.}
mediram, em um dos MEMS avaliados, sensibilidade entre \num{63} e
\SI{111}{\milli\volt\per\gram}, contra um piezo de referência de
\SI{100}{\milli\volt\per\gram}~\cite{albarbar2008}. O registro de auditoria contém apenas
esses dois valores; a sensibilidade nominal de \SIrange{450}{550}{\milli\volt\per\gram},
com que a faixa medida costuma ser comparada, vem das notas de leitura do projeto como
valor de folha de dados do fabricante e não do artigo. Se esse par estiver correto, a
discrepância vai de cerca de quatro a nove vezes ($450/111 \approx 4,1$ e $550/63 \approx
8,7$); como a rotulagem dos sensores no artigo é internamente inconsistente, o dado é
usado aqui apenas como ilustração qualitativa da dispersão de fator de escala em MEMS não
calibrados, e a mesma ressalva vale no Eixo~3.

No polo oposto, Landi \emph{et al.} descrevem o sensor tido como apto: calibração de
\SIrange{10}{25000}{\hertz} conforme ISO 16063-21:2003, sensibilidade a \SI{10}{\hertz} de
\SI{2,004}{\milli\volt} por \si{\meter\per\second\squared} e banda de exatidão de
\SI{10}{\percent} até \SI{12}{\kilo\hertz}~\cite{landi2023}. Rodrigues \emph{et al.}
dão apoio modesto ao MPU6050, mais uniforme que o ADXL345 sob o mesmo sinal aleatório de
baixa intensidade, com a ressalva dos autores de que explicar as discrepâncias exigiria
mais lotes e fabricantes, e sem calibração rastreável~\cite{rodrigues2021}. Um valor em
\si{\milli\meter\per\second} vindo de um MPU6050 não calibrado herda, portanto, fator de
escala não verificado, e compará-lo à fronteira de \SI{1,4}{\milli\meter\per\second} é
exercício de aparência normativa. Isso vale mesmo quando a máquina medida estiver dentro
do escopo da norma, o que na validação em MAFAULDA não é o caso.

\subsubsection{O orçamento de ruído, que é a objeção mais dura}

Falta enfrentar o número que decide a questão antes mesmo da banda, e que estava sendo
omitido enquanto três parágrafos discutiam calibração. Literatura técnica da Analog
Devices estabelece que resolver a fronteira da zona ``boa'' de uma máquina de
Classe~2 ($v_{\min} = \SI{1,12}{\milli\meter\per\second}$) a \SI{10}{\hertz} exige
ruído de aceleração abaixo de \num{7,18}~\mgf. Esse orçamento circula fora
da revisão por pares, em material de fabricante; a auditoria não conseguiu localizá-lo no
artigo de Looney, que é a referência habitualmente citada para ele, e sua
atribuição precisa permanece pendente~\cite{looney2017}. Com essa ressalva de
proveniência, a conta própria é direta: \num{400}~\dens{} --- valor de folha de dados do
MPU6050, não verificado experimentalmente --- integrados sobre a banda de
\SI{990}{\hertz} dão $400\times10^{-6}\sqrt{990} \approx \num{12,6}$~\mgf, cerca de
\num{1,8} vezes acima do orçamento.

A conclusão fortalece a tese deste eixo em vez de enfraquecê-la: a limitação é
\emph{dupla}. Mesmo que a banda fosse suficiente, o piso de ruído do sensor não resolveria
a fronteira A/B da norma; e mesmo que o piso de ruído fosse suficiente, a banda cobriria
metade do exigido. Rastrear o documento correto da Analog Devices, e refazer a conta com
uma densidade de ruído medida em vez de declarada, são as duas providências que faltam
para que esse argumento passe de estimativa a resultado.

Resta o vazio. O par de baixo custo mais próximo em hardware, Kolok \emph{et al.}, combina
MPU6050, ESP32 e Isolation Forest e reporta \SI{72,9}{\percent} $\pm$~\SI{2,4}{\percent}
de acurácia, com revocação de \num{0,679} e ROC-AUC de \num{0,780}; o trabalho encerra em
valor eficaz de aceleração e FFT, sem integração para velocidade, sem critério em
\si{\milli\meter\per\second} e sem menção localizável à ISO 10816 ou
20816~\cite{kolok2025}. A pergunta específica segue sem resposta publicada: o que uma IMU
de consumo com taxa de saída de \SI{1}{\kilo\hertz} pode e não pode afirmar sobre
severidade normativa. É nessa pergunta que este trabalho se situa.

\subsection{Eixo 5 --- Sistemas IoT comparáveis, invólucro e restrições térmicas}

A literatura de nós sem fio para manutenção preditiva divide-se em duas populações que
raramente se citam. De um lado, produtos comerciais, que compartilham três atributos:
banda de medição uma ordem de grandeza acima da do Kaelix, caminho mecânico rígido até o
elemento sensor, obtido por base metálica aparafusada ou por base colada ao ativo, e
alimentação por célula primária. Convém não generalizar mais do que isso: apenas um dos
três usa elemento piezoelétrico, e dois têm corpo declaradamente plástico. Do outro lado,
nós acadêmicos de baixo custo, que convergem para a mesma arquitetura do Kaelix: ESP32,
IMU MEMS na ordem de \SI{1}{\kilo\hertz} e detecção não supervisionada por \emph{Isolation
Forest}. Lê-las lado a lado é o que separa, no projeto, escolha defensável de limitação
herdada da classe de hardware.

\subsubsection{Nós de baixo custo: o que a literatura de fato autoriza concluir}

O trabalho de Kolok \emph{et al.} é o análogo mais próximo em hardware e em modo de
falha: ESP32-C6 com MPU6050 e \emph{Isolation Forest} de \num{100} árvores com
\texttt{contamination='auto'}, abaixo de 30~EUR de hardware e \SI{300}{\milli\watt},
reportando acurácia de \num{0,729}, F1 de \num{0,715} e ROC-AUC de \num{0,780}, com
validação cruzada de cinco dobras em \SI{72,9}{\percent} $\pm$~\SI{2,4}{\percent} sobre
cerca de \num{900} amostras e \emph{um único} modo de falha~\cite{kolok2025}. O dado valida
a arquitetura como publicável. O que ele não faz é fixar teto de desempenho da classe, e
dizer que fixa seria erro triplo: a amostra é de um único estudo com um único modo de
falha; o \emph{streaming} efetivo é de uma linha a cada \SI{0,5}{\second}, cerca de
\SI{2}{\hertz}, contra taxa interna de \SI{1}{\kilo\hertz}~\cite{kolok2025}; e as notas de
leitura do projeto registram que o filtro passa-baixas interno do MPU6050 foi configurado
em \SI{5}{\hertz}, dado ausente do registro de auditoria e que precisa ser reconfirmado.
Se isso se confirmar, aquele nó operou com banda útil duas ordens de grandeza abaixo da do
Kaelix, e o que ele demonstra é que a arquitetura ESP32 + MPU6050 + Isolation Forest é
publicável para desbalanceamento em banda baixíssima, sem informação sobre rolamento e sem
informação sobre a banda deste projeto.

Chevtchenko \emph{et al.} são a comparação mais direta, porque o experimento é embarcado e
mede três detectores no mesmo nó e no mesmo conjunto: ESP32 com IMU ISM330DLCTR a
\SI{1}{\kilo\hertz}, dois microfones SPH0645LM4H-B a \SI{20}{\kilo\hertz}, fixação por ímã
de neodímio, sobre \num{3750} amostras normais e \num{2250} anômalas de três motores de
indução de 1~hp a \num{1000}~rpm. Os resultados de validação e teste, em sensibilidade e
especificidade, são: Isolation Forest \SI{88,9}{\percent}/\SI{86,3}{\percent} e
\SI{61,6}{\percent}/\SI{67,4}{\percent}, com inferência de \SI{21,40}{\milli\second}; LOF
\SI{90,8}{\percent}/\SI{77,6}{\percent} e \SI{74,5}{\percent}/\SI{72,1}{\percent}, com
\SI{0,81}{\milli\second}; \emph{one-class} SVM \SI{73,1}{\percent}/\SI{47,9}{\percent} e
\SI{63,4}{\percent}/\SI{36,5}{\percent}, com
\SI{0,43}{\milli\second}~\cite{chevtchenko2023}. No mesmo hardware e no mesmo dado, o LOF
domina o Isolation Forest tanto em falso alarme, com \SI{74,5}{\percent} contra
\SI{61,6}{\percent} de especificidade em validação, quanto em latência, sendo cerca de
\num{26} vezes mais rápido. A objeção mais direta ao desenho do Kaelix está, portanto,
dentro de uma fonte que o projeto já cita, e não em outro conjunto de dados.

A resposta declarada é de memória, não de acurácia, e precisa ser verificada: o LOF é um
método baseado em vizinhança, e sua inferência exige manter o conjunto de referência ---
as próprias amostras de treino --- e calcular distâncias aos $k$ vizinhos a cada janela,
de modo que o custo de memória cresce com o tamanho do conjunto de treino; a floresta de
isolamento armazena apenas limiares e índices, com custo fixo depois de treinada. Em um
ESP32-S3 com folga de RAM medida em centenas de kilobytes, essa diferença pode decidir a
questão --- ou não decidir, e nesse caso o LOF é a escolha melhor. Como o orçamento
embarcado do Kaelix ainda não foi medido, a permanência no Isolation Forest é, hoje,
hipótese de projeto e não conclusão.

O mesmo trabalho traz um segundo achado que é restrição de implantação e não detalhe: os
autores observam que o terceiro motor, com cerca de um ano de uso, tem assinatura normal
distinta dos motores novos~\cite{chevtchenko2023}. Um modelo único com limiar fixo herda
essa deriva do que conta como ``normal'' entre máquinas e ao longo da vida, o que é o
cenário em que a calibração por risco alvo discutida no Eixo~1 deixa de ser refinamento e
passa a ser requisito.

A taxa de falso alarme de \SIrange{32,6}{38,4}{\percent} medida ali é da mesma ordem de
grandeza do achado próprio de que \texttt{contamination='auto'} marcaria cerca de
\SI{42}{\percent} da operação normal do Kaelix como anômala. É preciso dizer o que essa
proximidade não autoriza: o registro verificado de Chevtchenko \emph{et al.} não declara
a política de limiar usada, e sem ela a comparação é de ordem de grandeza entre conjuntos,
máquinas e cortes diferentes. Não é correlato empírico controlado. O Eixo~1 e a
Tabela~\ref{tab:fpr} tratam o ponto.

Há ainda uma objeção de viabilidade computacional, e ela precisa ser enunciada com o
regime de supervisão à vista. Pipelines mais pesados cabem em microcontroladores
comparáveis: \SI{94,1}{\percent} de acurácia diagnóstica com descritores de vibração e
\SI{95,5}{\percent} com componentes principais de corrente em RP2040, contra
\SI{83,2}{\percent} com estatísticas manuais, com destilação reduzindo a pegada de
\SI{0,42}{\mega\byte} para \SI{0,15}{\mega\byte}~\cite{uselis2025}; e uma rede quantizada
em INT8 no mesmo ESP32-S3 com latência média determinística de
\SI{42,3}{\milli\second}~\cite{kumar2026}. Esses números \emph{não} são comparáveis ao do
Kaelix: são classificadores supervisionados, treinados com rótulos, sobre outros conjuntos
de dados, enquanto o Kaelix aposta em detecção não supervisionada. O que eles autorizam
concluir é apenas o que importa aqui: o orçamento de um microcontrolador desta classe
comporta \emph{pipelines} bem mais caros que uma floresta de isolamento sobre quatro
descritores, de modo que o custo computacional não é, por si só, justificativa para a
escolha do detector.

\subsubsection{Banda, invólucro e caminho mecânico}

É na física do dispositivo que a literatura contraria o projeto mais diretamente. Os
produtos comerciais resolvem justamente a banda que ele descarta, e a
Tabela~\ref{tab:comerciais} os apresenta com as grandezas separadas, banda de exatidão de
um lado e taxa de amostragem de outro, porque misturá-las é o erro que este documento
cometia ao descrever o concorrente com o rigor de Nyquist que aplicava a si mesmo.

\begin{table}[!htb]
\centering
\caption{\textbf{Sensores comerciais de vibração sem fio, com as grandezas separadas.}
Banda de exatidão e taxa de amostragem são declaradas em colunas distintas; onde o
datasheet informa apenas taxa, a frequência máxima analisável é $f_s/2$. Todos os valores
vêm das folhas de dados citadas.}
\label{tab:comerciais}
\footnotesize
\setlength{\tabcolsep}{4pt}
\begin{tabular}{@{}p{2.2cm}p{1.7cm}p{4.9cm}p{2.6cm}p{2.4cm}@{}}
\toprule
\textbf{Produto} & \textbf{Elemento} & \textbf{Banda de aceleração} & \textbf{Fixação} & \textbf{Bateria} \\
\midrule
Ranger Pro \cite{bentlynevada2025}
  & piezo cerâmico
  & Z \SIrange{5}{10000}{\hertz} ($\pm$\SI{3}{\decibel}); X/Y \SIrange{5}{4000}{\hertz}. Velocidade \SIrange{5}{2000}{\hertz}
  & rosca M6$\times$1 (\emph{stud}), sem ímã
  & LiSOCl$_2$ D, substituível \\[3pt]
Phantom Gen~3 \cite{erbessd2025}
  & MEMS digital
  & V11E: X/Y \SIrange{0,5}{10}{\kilo\hertz}, Z \SIrange{0,5}{5,1}{\kilo\hertz}; $f_s$ até \SI{25,6}{\kilo\hertz}. V10E: X/Y \SIrange{0,5}{4}{\kilo\hertz}
  & base em aço inox, parafuso 1/4-28~UNF
  & CR2477 primária, não substituível \\[3pt]
Smart Trac \cite{tractian2022}
  & MEMS digital
  & faixa declarada \SIrange{1}{32000}{\hertz}, não separada de $f_s$ no datasheet (se for $f_s$, $F_{\max}=\SI{16}{\kilo\hertz}$)
  & ímã, rosca ou epóxi; base idealmente colada
  & lítio primária, não substituível \\[3pt]
Kaelix (projeto)
  & MEMS (MPU6050)
  & ODR \SI{1}{\kilo\hertz} $\Rightarrow$ $F_{\max}=\SI{500}{\hertz}$, sem DLPF configurado
  & ímã
  & LiPo recarregável \\
\bottomrule
\end{tabular}
\end{table}

O mecanismo que torna essa diferença decisiva já foi estabelecido no Eixo~3: o defeito
localizado é detectado pela ressonância estrutural de alta frequência que os impactos
excitam, e a taxa de repetição aparece como modulação dessa portadora, o que faz da
demodulação de envelope numa banda alta o método de referência~\cite{randall2011}, e da
curtose uma função da frequência cuja banda ótima precisa ser escolhida~\cite{antoni2006}.
Os três produtos, \emph{ao menos em um eixo}, medem uma ordem de grandeza acima do teto do
MPU6050 e alcançam essa banda de ressonância; a assimetria importa e estava sendo omitida:
no Ranger Pro os \SI{10}{\kilo\hertz} existem apenas no eixo Z, com X e Y limitados a
\SI{4}{\kilo\hertz}~\cite{bentlynevada2025}, e no Phantom a variante de alta sensibilidade
EPH-V10E cobre apenas \SIrange{0,5}{4}{\kilo\hertz} em X e Y e
\SIrange{0,5}{1,8}{\kilo\hertz} em Z~\cite{erbessd2025}. Dois deles entregam também
velocidade: o Ranger Pro mede velocidade de \SIrange{5}{2000}{\hertz}, cobrindo
integralmente a banda normativa de \SIrange{10}{1000}{\hertz}~\cite{bentlynevada2025}, e o
Smart Trac reporta velocidade e aceleração com valor eficaz, pico, pico a pico e fator de
crista, além de ferramentas de BPFO, BPFI, BSF, FTF e envelope~\cite{tractian2022}.

A suposição de que carcaça plástica seria inviável não se sustenta: o Smart Trac declara
corpo externo em termoplástico policarbonato LEXAN com encapsulamento eletrônico integral
em resina, IP69K e operação de \SIrange{-10}{90}{\celsius}~\cite{tractian2022}, e o Phantom
declara ``\textit{Housing material: ABS, stainless steel}''~\cite{erbessd2025}. Uma nota de
leitura anterior do projeto descrevia o corpo do Smart Trac como metálico; a conferência
da folha de dados primária resolveu a divergência a favor do policarbonato.
Mas a conclusão que interessa ao Kaelix é outra, e vai na direção contrária: em nenhum dos
três o \emph{caminho mecânico} entre o ativo e o elemento sensor é plástico ou magnético.
O Ranger Pro tem corpo em aço inoxidável 316 com topo em PPS reforçado e fixa por rosca
M6$\times$1, sem ímã~\cite{bentlynevada2025}; o Phantom tem base em aço inoxidável com
parafuso 1/4-28~UNF~\cite{erbessd2025}; e o Smart Trac, embora aceite ímã, rosca ou epóxi,
declara que ``\textit{the sensor base has a magnet to help with the installation process.
Ideally, the base should be glued to the asset}'', com cola bicomponente acompanhando o
produto e porca de travamento~\cite{tractian2022}. Nenhum dos três aceita o ímã como
caminho mecânico definitivo. O nó \emph{open hardware} mais próximo do Kaelix vai no mesmo
sentido: usina pé e suporte de PCB em alumínio, deixando só a tampa impressa, porque o
alumínio é preferido por estabilidade de temperatura e maior rigidez para transmitir
vibração aos sensores~\cite{jakobsen2024}.

\subsubsection{A frequência de montagem: estimativa, não medição}

O achado próprio de que o invólucro em ASA forma sistema massa-mola com frequência natural
de \SI{557}{\hertz}, dentro da banda de medição, e transmissibilidade da ordem de cinco
vezes em \SI{500}{\hertz}, precisa ser apresentado com duas ressalvas que estavam
faltando. A primeira é de natureza: são saída de um modelo analítico de um grau de
liberdade, com rigidez em série estimada a partir das cotas do desenho e coeficiente de
amortecimento \emph{arbitrado} em $\zeta = \num{0,03}$. Esse valor é plausível para peça
impressa em ASA, jamais foi medido, e a transmissibilidade de pico depende dele
diretamente. O resultado é ainda sensível a um parâmetro geométrico incerto: dobrando a
espessura da base, a frequência sobe para \SI{750}{\hertz}. São, portanto, estimativas com
faixa de incerteza, e não medições. Medir a frequência natural do conjunto montado por
excitação de impacto é ensaio de bancada barato, e é o que falta.

A segunda ressalva é que ``não há precedente publicado'' era afirmação forte demais. Existe
critério consolidado em acelerometria: a frequência de montagem deve ser pelo menos três
vezes a maior frequência de interesse, o que corresponderia a \SI{3}{\kilo\hertz} para a
banda normativa. O levantamento do projeto o associa à ISO 5348, que também tabula a
degradação por método de fixação e seria a referência óbvia para uma fixação magnética.
Essa norma não foi conferida em fonte primária nesta auditoria e não integra a
bibliografia, mas sua existência muda o estatuto do achado: os \SI{557}{\hertz} não são um
resultado sem termo de comparação, são a reprovação previsível em um critério conhecido,
por um fator de mais de cinco. A ressonância própria da fixação magnética, que degrada
ainda mais o caminho, permanece ausente da análise.

O que continua sem precedente é o número publicado. Nenhum dos nós acadêmicos revisados
reporta a frequência natural do conjunto montado: Jakobsen registra apenas
autorressonâncias de PCB e estrutura entre \SIrange{6}{9}{\kilo\hertz}~\cite{jakobsen2024},
Antonini \emph{et al.} adiam o invólucro para trabalho futuro~\cite{antonini2023}, e
não localizamos, na parte acessível do trabalho de Bisio \emph{et al.}, publicação da
frequência natural do conjunto montado. O texto integral está atrás de assinatura, e
essa é a fronteira da nossa busca, não uma afirmação sobre o conteúdo do
artigo~\cite{bisio2025}.

\subsubsection{Bateria e margem térmica}

Na energia, a literatura contraria o projeto com clareza. Nabavi \emph{et al.} recomendam
explicitamente ``Li-SOCl$_2$ batteries for low-power, long-term deployments in extreme
environments'' e reservam as químicas recarregáveis a ``high-power applications requiring
frequent cycling''~\cite{nabavi2025}; a prática confirma em todos os sistemas revisados:
LS14500 com supercapacitor de \SI{100}{\milli\farad} para o pulso de
\SI{120}{\milli\ampere} do LoRa~\cite{jakobsen2024}, célula LiSOCl$_2$ D substituível com
até cinco anos~\cite{bentlynevada2025}, CR2477 para \num{65000}
medições~\cite{erbessd2025} e lítio não substituível por três anos~\cite{tractian2022}. A
opção pela LiPo recarregável é, nesse ponto, contrariada pela literatura e deve ser
tratada como pergunta de projeto, não como resultado.

Já a \emph{forma} do limite térmico encontra apoio: no Ranger Pro o aparelho opera de
\SIrange{-40}{85}{\celsius}, mas a temperatura ambiente admissível em área classificada é
fixada \emph{pela bateria}, em \SI{80}{\celsius} ou \SI{70}{\celsius} conforme o
modelo~\cite{bentlynevada2025}. É o mesmo padrão do Kaelix, cuja margem é imposta pela
célula e não pelo ASA. O valor específico, porém, não vem da literatura de manutenção
preditiva, e é preciso dizer de onde ele vem: o limite de \SI{45}{\celsius} para carga é
especificação de folha de dados da célula e requisito dos regimes de carga das normas de
segurança de baterias de lítio, documentos que não foram conferidos nesta auditoria e não
constam da bibliografia. A revisão sobre efeitos térmicos em íon-lítio aqui citada
estabelece que a temperatura ``significantly impacts on the performance of lithium-ion
batteries and also limits the application''~\cite{ma2018}, mas o resumo não menciona limite
de temperatura de carga e o texto integral está atrás de assinatura. Ela ampara, portanto,
o mecanismo geral, e não o número.

A lacuna, portanto, é mais estreita do que se afirmava. Não é verdade que o assunto esteja
``ausente da conversa'': a razão física do limite de carga é amplamente publicada na
literatura de envelhecimento de íon-lítio. O que não se localizou foi orçamento térmico
publicado para carga de célula secundária dentro de invólucro montado em carcaça quente de
motor. A razão é estrutural: o campo usa célula primária e, portanto, nunca carrega nada.

\section{Síntese}

Lidos em conjunto, os quatro eixos não formam uma lista de pontos fortes e fracos: formam
uma \emph{cadeia com ordem de precedência}. Primeiro vem a aquisição: piso de ruído, banda
e caminho mecânico. Nenhuma decisão posterior recupera informação que o sensor não
coletou. Depois vem a estimação, que decide o que se faz com a informação disponível:
limiar, métrica e protocolo de validação. Por último, a implantação --- energia, térmica e
invólucro --- decide por quanto tempo o dispositivo continua fazendo isso. A conclusão
desta revisão é que a literatura é mais forte onde o Kaelix é mais frágil, o primeiro elo,
e mais frágil onde o Kaelix tem algo a dizer, o segundo.

\subsection{O que a literatura sustenta}

O problema do limiar é real, e não é preciosismo. O escore do Isolation Forest tem
\num{0,5} como referência interpretativa e não como corte calibrado~\cite{liu2008}, e a
implementação de referência apenas transporta essa interpretação para um \emph{offset}
fixo, deixando o ranqueamento intocado~\cite{scikitlearn2026}. Os dois trabalhos que
reportam bom desempenho com o detector abandonam o padrão: um fixa o escore normalizado
em \num{0,75}~\cite{antonini2023}, o outro deriva o corte da contaminação
conhecida~\cite{brito2022}. Existe ainda alternativa principiada, com risco alvo em vez de
fração de anomalias adivinhada~\cite{siffer2017}. A decisão de reportar pAUC em faixa de
baixa FPR, e não apenas AUC global, segue o benchmark de referência para o mesmo regime de
treino~\cite{koizumi2020} e encontra eco em quem reporta PR-AUC~\cite{brito2022}; o
resultado próprio de que a pAUC cai \SI{16,2}{\percent} onde a AUC cai \SI{3,7}{\percent}
--- em sinal sintetizado, entre as duas cadeias de aquisição --- mostra por que a escolha
importa.

A viabilidade computacional está estabelecida: o Isolation Forest treina e executa dentro
de um ESP32~\cite{antonini2023}, existe variante desenhada para reduzir seu custo em
TinyML~\cite{barbariol2022}, e microcontroladores da mesma classe comportam
\emph{pipelines} consideravelmente mais caros~\cite{uselis2025,kumar2026}. A arquitetura
ESP32 + MPU6050 + detecção não supervisionada é publicável~\cite{kolok2025}. A exigência
de integrar aceleração para velocidade é normativa~\cite{iso208163_2022} e o método
correto está estabelecido: frequência, não trapézio no tempo~\cite{brandt2014}. E a forma
do limite térmico do Kaelix, margem imposta pela célula e não pelo invólucro, tem
precedente comercial direto~\cite{bentlynevada2025}.

Há, por fim, uma convergência que o projeto deveria reivindicar explicitamente. Os
resultados próprios em dados reais, a \SI{2}{\percent} de falso alarme, dão
\SI{75,0}{\percent} em desbalanceamento e entre \num{15} e \SI{21}{\percent} em defeito de
rolamento puro. Eles reproduzem, por medição, a divisão de trabalho que Jakobsen
implementou por projeto ao reservar um acelerômetro de \SI{400}{\hertz} para
desbalanceamento e outro de \SI{40}{\kilo\hertz} para rolamento~\cite{jakobsen2024}, e o
padrão que Kolok \emph{et al.} obtiveram ao demonstrar desbalanceamento e adiar
rolamento~\cite{kolok2025}. Três caminhos independentes chegam à mesma fronteira.

\subsection{O que a literatura contesta}

As objeções são de peso, e nenhuma delas se resolve com ajuste de \emph{software}.

\textbf{A banda.} Nenhum estudo de adequação revisado testou sensor com limite superior
abaixo de \SI{1500}{\hertz}~\cite{albarbar2008}, os grupos que querem diagnosticar
rolamento com MEMS escolhem sensores de dezenas de
quilohertz~\cite{landi2023,jakobsen2024,sharma2024}, e os três produtos comerciais medem,
ao menos em um eixo, uma ordem de grandeza acima do teto do
MPU6050~\cite{bentlynevada2025,erbessd2025,tractian2022}. O mecanismo é conhecido e não
admite contorno: a assinatura do defeito viaja numa portadora de alta
frequência~\cite{randall2011,antoni2006,antoni2007}.

\textbf{O piso de ruído.} A objeção é anterior à banda, e o projeto a possuía sem tê-la
usado: com \num{400}~\dens{} de folha de dados integrados sobre a banda normativa, o
sensor fica cerca de \num{1,8} vezes acima do ruído admissível para resolver a fronteira
A/B em \SI{10}{\hertz}. A limitação é dupla, e a segunda metade dela não depende de
nenhuma escolha de processamento.

\textbf{O detector.} A maior comparação restrita a métodos não supervisionados recomenda o
Extended Isolation Forest~\cite{bouman2024}; com \SI{1}{\percent} de rótulos, métodos
semissupervisionados superam o melhor não supervisionado~\cite{han2022}; e, no mesmo
hardware embarcado e no mesmo conjunto de dados, o LOF é mais específico e cerca de
\num{26} vezes mais rápido que o Isolation Forest~\cite{chevtchenko2023}. Nenhuma dessas
três objeções foi respondida por medição neste projeto; as respostas oferecidas no Eixo~1 e
no Eixo~5 são critérios de projeto declarados, e permanecem hipóteses até que o orçamento
embarcado seja medido.

\textbf{O falso alarme.} Três implementações independentes de Isolation Forest em nó
embarcado ficam entre \num{22} e \SI{38}{\percent} de falso alarme sobre operação
normal~\cite{kolok2025,chevtchenko2023}. O achado próprio de \SI{42}{\percent} com
\texttt{contamination='auto'} não é anomalia deste projeto: é o comportamento esperado do
parâmetro, e a literatura já o exibia sem nomeá-lo.

\textbf{O caminho mecânico.} Nenhum dos três produtos comerciais aceita o ímã como fixação
definitiva~\cite{bentlynevada2025,erbessd2025,tractian2022}, e o nó acadêmico mais próximo
usina alumínio precisamente por rigidez~\cite{jakobsen2024}. A frequência de montagem
estimada em \SI{557}{\hertz} não é uma curiosidade: é reprovação, por fator superior a
cinco, em um critério de acelerometria consolidado.

\textbf{A energia.} A literatura de seleção de bateria para IoT recomenda célula primária
para implantação de longa duração e reserva as recarregáveis a aplicações de ciclagem
frequente~\cite{nabavi2025}, e todos os sistemas revisados seguem essa
recomendação~\cite{jakobsen2024,bentlynevada2025,erbessd2025,tractian2022}. A LiPo
recarregável do Kaelix é a escolha isolada.

\subsection{O que se conclui}

Duas consequências, e uma observação que a literatura não explica.

A primeira consequência é de posicionamento, não de engenharia. Com a banda e o piso de
ruído que tem, o Kaelix não é um diagnosticador de rolamento, e insistir em apresentá-lo
como tal é o caminho mais curto para uma banca desmontá-lo com um único gráfico. O que ele
é, e o que os resultados próprios sustentam, é um detector de desbalanceamento e de
desvios de operação com banda declarada, um canal de telemetria de baixo custo e uma
plataforma sobre a qual questões de limiar e de métrica podem ser estudadas honestamente.
Reivindicar isso é mais forte do que reivindicar o resto e ter de recuar.

A segunda é de método. As contribuições possíveis deste trabalho estão quase todas no
segundo elo da cadeia: limiar calibrado por risco alvo, pAUC no ponto de operação,
validação por ensaio, estratificação por confundidor. São contribuições reais porque
a literatura de vibração embarcada as ignora sistematicamente. Mas elas só valem se o
primeiro elo estiver declarado, e não disfarçado. Um resultado de \SI{75}{\percent} em
desbalanceamento com banda declarada e falso alarme fixado em \SI{2}{\percent} é
publicável; o mesmo número apresentado como ``detecção de falhas em motores'' não é.

A observação que a literatura não explica é esta: dos três modos de falha medidos em dados
reais, o desalinhamento é o pior, com \SIrange{5}{7}{\percent} de detecção a
\SI{2}{\percent} de falso alarme, e ele \emph{não} é um problema de banda. Desalinhamento
se manifesta em harmônicas da rotação, muito abaixo de \SI{500}{\hertz}, plenamente dentro
da banda do sensor. Se o desempenho é ruim ali, a causa está em outro lugar: nos
descritores escolhidos. Valor eficaz, curtose, fator de crista e frequência dominante
descrevem energia e impulsividade e não capturam a estrutura harmônica e de fase em que o
desalinhamento vive. Nenhum dos trabalhos revisados oferece explicação para esse caso,
porque nenhum deles mede desalinhamento com detector não supervisionado em banda limitada.
É a pista mais direta do conjunto, e é acionável: acrescentar descritores de razão
harmônica não custa banda, não custa sensor e não custa dinheiro.

\section{Lacunas e oportunidades}

Cada lacuna abaixo declara o que foi procurado, onde, e o que não foi encontrado. Duas
lacunas anteriormente enunciadas caíram na auditoria e estão registradas aqui como caídas,
porque uma revisão que só apresenta as lacunas que sobreviveram esconde o método.

\paragraph{Caiu: ``não existe FPR publicado para Isolation Forest embarcado''.} Existe.
Chevtchenko \emph{et al.} medem especificidade de \SI{61,6}{\percent} e
\SI{67,4}{\percent} para Isolation Forest em nó ESP32~\cite{chevtchenko2023}, e o valor de
Kolok \emph{et al.} é derivável de precisão, revocação e acurácia
publicadas~\cite{kolok2025}. A Tabela~\ref{tab:fpr} os reúne.

\paragraph{Caiu: ``ninguém quantifica a fração de normais marcada por
\texttt{contamination='auto'}''.} Kolok \emph{et al.} usam exatamente essa configuração
sobre vibração de MPU6050, e a fração é derivável dos números publicados.

\paragraph{Lacuna 1 --- efeito isolado do parâmetro de contaminação.} Procurou-se, entre os
trabalhos que executam Isolation Forest em microcontrolador e entre os comparativos de
detectores, algum que variasse \texttt{contamination} sobre a \emph{mesma} base e o
\emph{mesmo} modelo e reportasse a fração de operação normal marcada. Não há:
Kolok \emph{et al.} usam \texttt{'auto'} sem variá-lo, Chevtchenko \emph{et al.} não
declaram a política de limiar no registro verificado, e Antonini \emph{et al.} fixam
\num{0,75} e declaram a avaliação de acurácia fora de escopo. O experimento que falta é
barato e definitivo: uma varredura do parâmetro contra FPR medido, em vibração industrial
de banda limitada.

\paragraph{Lacuna 2 --- falso alarme de detector treinado no próprio dispositivo.} O único
trabalho revisado que treina \emph{on-device} é justamente o que não avalia
acurácia~\cite{antonini2023}. Não se localizou nenhum FPR reportado para Isolation Forest
treinado no microcontrolador, o que é exatamente o modo de operação que um dispositivo sem
histórico rotulado precisa adotar.

\paragraph{Lacuna 3 --- contraste diagnóstico como função contínua da banda.} A literatura
discute o efeito da banda qualitativamente, em termos como ``estágio inicial \emph{versus}
avançado'' e ``\SI{2,5}{\kilo\hertz} \emph{versus} \SI{10}{\kilo\hertz}'', e o único
estudo revisado que faz varredura de taxa de amostragem o faz para acurácia de
classificador supervisionado, não para contraste do indicador~\cite{elbouharrouti2024}. A
oportunidade é executar a varredura com \emph{método fixo} sobre a mesma gravação decimada
em vários $f_s$, e reportar a decomposição entre efeito de banda e efeito de perda do
envelope, o que o par \num{413}$\times$/\num{20}$\times$ deste projeto ainda não separa.

\paragraph{Lacuna 4 --- artefato de aliasing sobre descritores estatísticos.} Procurou-se,
na literatura revisada por pares de diagnóstico de rolamento, descrição experimental do
efeito de ausência de filtro anti-\emph{aliasing} sobre curtose e fator de crista em IMU de
consumo. Não foi encontrada; o efeito aparece apenas em notas de aplicação de fabricante,
não conferidas nesta auditoria. É um resultado publicável e acionável: medir
os mesmos descritores com e sem DLPF configurado, no mesmo sinal.

\paragraph{Lacuna 5 --- custo do truncamento sobre o indicador normativo.}
Brandt e Brincker quantificam o erro dos métodos de integração~\cite{brandt2014}, mas
não se localizou a aplicação desses métodos ao indicador da norma. Mais importante: não se
localizou publicação de quanto um dispositivo truncado em \SI{500}{\hertz} subestima
$v_{\mathrm{rms}}$ frente à banda normativa de \SIrange{10}{1000}{\hertz}. É calculável
com os dados já disponíveis, e transforma retórica em número.

\paragraph{Lacuna 6 --- detectabilidade limitada por piso de ruído.} Não se localizou, na
literatura revisada, análise de defeito mínimo detectável em função da densidade de ruído
para IMUs de consumo. O orçamento que existe é de literatura de fabricante e sua atribuição
não pôde ser confirmada~\cite{looney2017}. Fechar essa conta com densidade de ruído medida,
e não declarada, é contribuição de instrumentação, e complementa a crítica de
ausência de calibração rastreável que os próprios estudos de MEMS
levantam~\cite{albarbar2008,rodrigues2021}.

\paragraph{Lacuna 7 --- frequência natural do conjunto montado.} Nenhum dos nós de baixo
custo revisados publica a frequência natural do conjunto invólucro-sensor-fixação:
Jakobsen reporta apenas autorressonâncias de PCB e estrutura entre
\SIrange{6}{9}{\kilo\hertz}~\cite{jakobsen2024}, Antonini \emph{et al.} adiam o
invólucro~\cite{antonini2023}, e não localizamos o dado na parte acessível de
Bisio \emph{et al.}~\cite{bisio2025}. A oportunidade não é a estimativa analítica, que
já existe e depende de um amortecimento arbitrado: é a \emph{medição} por excitação de
impacto do conjunto montado, incluindo a ressonância própria da fixação magnética, e o
confronto com o critério de acelerometria.

\paragraph{Lacuna 8 --- orçamento térmico de carga de célula secundária.} Procurou-se
literatura sobre carga de bateria recarregável em nó de monitoramento montado em carcaça
de motor. O campo usa célula primária e portanto nunca carrega nada~\cite{nabavi2025}; a
revisão térmica disponível estabelece o mecanismo mas não o limite de
carga~\cite{ma2018}. O limite de \SI{45}{\celsius} vem da folha de dados da célula e das
normas de segurança aplicáveis, não da literatura de manutenção preditiva. O que falta
publicado é o orçamento térmico completo: carcaça quente, invólucro fechado, célula
carregando. É ele que decide se a escolha de bateria do projeto sobrevive.

\paragraph{Lacuna 9 --- separar efeito de banda de efeito de montagem.} Não se localizou
benchmark aberto com sinais gravados \emph{simultaneamente} por MEMS de banda limitada e
por piezoelétrico no mesmo ponto de medição. Sem isso, os efeitos de banda (Eixo~3) e de
caminho mecânico (Eixo~5) permanecem confundidos em qualquer medição de campo, inclusive
nas deste projeto.

\section{Referências}
\label{sec:bib}

\begin{thebibliography}{99}
\setlength{\itemsep}{2pt}

\bibitem{albarbar2008}
A.~Albarbar, S.~Mekid, A.~Starr e R.~Pietruszkiewicz,
\newblock Suitability of MEMS accelerometers for condition monitoring: an experimental study,
\newblock \emph{Sensors}, v.~8, n.~2, p.~784--799, 2008.
\newblock DOI: 10.3390/s8020784.

\bibitem{antoni2006}
J.~Antoni e R.~B.~Randall,
\newblock The spectral kurtosis: application to the vibratory surveillance and diagnostics of rotating machines,
\newblock \emph{Mechanical Systems and Signal Processing}, v.~20, n.~2, p.~308--331, 2006.
\newblock DOI: 10.1016/j.ymssp.2004.09.002.

\bibitem{antoni2007}
J.~Antoni,
\newblock Fast computation of the kurtogram for the detection of transient faults,
\newblock \emph{Mechanical Systems and Signal Processing}, v.~21, n.~1, p.~108--124, 2007.
\newblock DOI: 10.1016/j.ymssp.2005.12.002.
\newblock Metadados conferidos no depósito do editor; resumo integral lido em fonte secundária.

\bibitem{antonini2023}
M.~Antonini, M.~Pincheira, M.~Vecchio e F.~Antonelli,
\newblock An adaptable and unsupervised TinyML anomaly detection system for extreme industrial environments,
\newblock \emph{Sensors}, v.~23, n.~4, art.~2344, 2023.
\newblock DOI: 10.3390/s23042344.

\bibitem{barbariol2022}
T.~Barbariol e G.~A.~Susto,
\newblock TiWS-iForest: isolation forest in weakly supervised and tiny ML scenarios,
\newblock \emph{Information Sciences}, v.~610, p.~126--143, 2022.
\newblock DOI: 10.1016/j.ins.2022.07.129. Preprint: arXiv:2111.15432.

\bibitem{bentlynevada2025}
Bently Nevada (Baker Hughes),
\newblock \emph{Ranger Pro wireless condition monitoring} --- folha de dados, doc.\ 125M5237 Rev.~U,
\newblock Baker Hughes, 16 jul.\ 2025. Revisão corrente Rev.~V (dez.\ 2025), com valores técnicos idênticos.
\newblock Literatura de fabricante, não revisada por pares.

\bibitem{bisio2025}
I.~Bisio, C.~Garibotto, A.~Grattarola, A.~Iscra, F.~Lavagetto, A.~Sciarrone e M.~Zerbino,
\newblock Design, implementation and performance of an IIoT node for vibration monitoring,
\newblock \emph{2025 IEEE International Symposium on Circuits and Systems (ISCAS)}, Londres, p.~1--4, 2025.
\newblock DOI: 10.1109/ISCAS56072.2025.11043851.
\newblock Texto integral não acessado (assinatura); citado apenas como contexto.

\bibitem{bouman2024}
R.~Bouman, Z.~Bukhsh e T.~Heskes,
\newblock Unsupervised anomaly detection algorithms on real-world data: how many do we need?,
\newblock \emph{Journal of Machine Learning Research}, v.~25, n.~105, p.~1--34, 2024.
\newblock JMLR não atribui DOI.

\bibitem{bouman2025}
R.~Bouman e T.~Heskes,
\newblock Autoencoders for anomaly detection are unreliable,
\newblock \emph{Preprint} arXiv (cs.LG), 23 jan.\ 2025.
\newblock arXiv:2501.13864. Sem revisão por pares até ago.\ 2026.

\bibitem{brandt2014}
A.~Brandt e R.~Brincker,
\newblock Integrating time signals in frequency domain: comparison with time domain integration,
\newblock \emph{Measurement}, v.~58, p.~511--519, 2014.
\newblock DOI: 10.1016/j.measurement.2014.09.004.

\bibitem{brito2022}
L.~C.~Brito, G.~A.~Susto, J.~N.~Brito e M.~A.~V.~Duarte,
\newblock An explainable artificial intelligence approach for unsupervised fault detection and diagnosis in rotating machinery,
\newblock \emph{Mechanical Systems and Signal Processing}, v.~163, art.~108105, 2022.
\newblock DOI: 10.1016/j.ymssp.2021.108105. Preprint: arXiv:2102.11848 (2021).

\bibitem{chevtchenko2023}
S.~F.~Chevtchenko, M.~C.~M.~dos Santos, D.~M.~Vieira, R.~L.~Mota, E.~Rocha, B.~V.~Cruz, D.~Araújo e E.~Andrade,
\newblock Predictive maintenance model based on anomaly detection in induction motors: a machine learning approach using real-time IoT data,
\newblock \emph{Proceedings of the 33rd European Safety and Reliability Conference (ESREL 2023)}, Research Publishing, Singapura, p.~3173--3180, 2023.
\newblock DOI: 10.3850/978-981-18-8071-1\_P578-cd.

\bibitem{elbouharrouti2024}
N.~El~Bouharrouti, D.~Morinigo-Sotelo e A.~Belahcen,
\newblock Multi-rate vibration signal analysis for bearing fault detection in induction machines using supervised learning classifiers,
\newblock \emph{Machines}, v.~12, n.~1, art.~17, 2024.
\newblock DOI: 10.3390/machines12010017. Publicado online em 27 dez.\ 2023.

\bibitem{erbessd2025}
Erbessd Instruments,
\newblock \emph{PHANTOM Gen~3 vibration sensor} --- folha de dados (EPH-V10E / EPH-V11E), rev.\ 11/2025,
\newblock Erbessd Instruments, 2025.
\newblock Literatura de fabricante, não revisada por pares.

\bibitem{fidali2024}
M.~Fidali, D.~Augustyn, J.~Ochmann e W.~Uchman,
\newblock Evaluation of the diagnostic sensitivity of digital vibration sensors based on capacitive MEMS accelerometers,
\newblock \emph{Sensors}, v.~24, n.~14, art.~4463, 2024.
\newblock DOI: 10.3390/s24144463.

\bibitem{han2022}
S.~Han, X.~Hu, H.~Huang, M.~Jiang e Y.~Zhao,
\newblock ADBench: anomaly detection benchmark,
\newblock \emph{Advances in Neural Information Processing Systems 35 (NeurIPS 2022)}, Datasets and Benchmarks Track, 2022.
\newblock Preprint: arXiv:2206.09426.

\bibitem{hermansa2022}
M.~Hermansa, M.~Kozielski, M.~Michalak, K.~Szczyrba, Ł.~Wróbel e M.~Sikora,
\newblock Sensor-based predictive maintenance with reduction of false alarms: a case study in heavy industry,
\newblock \emph{Sensors}, v.~22, n.~1, art.~226, 2022.
\newblock DOI: 10.3390/s22010226.

\bibitem{iso108163_2009}
ISO/TC 108/SC 2,
\newblock \emph{ISO 10816-3:2009 --- Mechanical vibration --- Evaluation of machine vibration by measurements on non-rotating parts --- Part 3: Industrial machines with nominal power above 15~kW and nominal speeds between 120~r/min and 15\,000~r/min when measured in situ},
\newblock 2.~ed., Genebra: ISO, 2009. Amd~1:2017.
\newblock \textbf{Norma retirada}, substituída pela ISO 20816-3:2022; mantida como referência por ser a que a prática industrial e parte da literatura ainda adotam.

\bibitem{iso208163_2022}
ISO/TC 108/SC 2,
\newblock \emph{ISO 20816-3:2022 --- Mechanical vibration --- Measurement and evaluation of machine vibration --- Part 3: Industrial machinery with a power rating above 15~kW and operating speeds between 120~r/min and 30\,000~r/min},
\newblock 1.~ed., Genebra: ISO, out.\ 2022.
\newblock Cancela e substitui a ISO 7919-3:2009, a ISO 10816-3:2009 e as respectivas emendas de 2017.

\bibitem{iso2954_2012}
ISO/TC 108/SC 3,
\newblock \emph{ISO 2954:2012 --- Mechanical vibration of rotating and reciprocating machinery --- Requirements for instruments for measuring vibration severity},
\newblock 2.~ed., Genebra: ISO, abr.\ 2012.
\newblock Cancela e substitui a ISO 2954:1975; vigência confirmada na revisão sistemática de 2023.

\bibitem{jakobsen2024}
M.~O.~Jakobsen,
\newblock Low cost MEMS accelerometer and microphone based condition monitoring sensor, with LoRa and Bluetooth Low Energy radio,
\newblock \emph{HardwareX}, v.~18, art.~e00525, 2024.
\newblock DOI: 10.1016/j.ohx.2024.e00525.

\bibitem{koizumi2020}
Y.~Koizumi, Y.~Kawaguchi, K.~Imoto, T.~Nakamura, Y.~Nikaido, R.~Tanabe, H.~Purohit, K.~Suefusa, T.~Endo, M.~Yasuda e N.~Harada,
\newblock Description and discussion on DCASE2020 challenge task 2: unsupervised anomalous sound detection for machine condition monitoring,
\newblock \emph{Proceedings of the DCASE2020 Workshop}, Tóquio, p.~81--85, 2020.
\newblock Preprint: arXiv:2006.05822.

\bibitem{kolok2025}
P.~Kolok, M.~Hodoň, P.~Ševčík, L.~Hotz e N.~Remy,
\newblock Low-cost IoT-based predictive maintenance using vibration,
\newblock \emph{Sensors}, v.~25, n.~21, art.~6610, 2025.
\newblock DOI: 10.3390/s25216610.
\newblock O artigo é internamente inconsistente: a discussão menciona \SI{93,4}{\percent} de acurácia, número ausente dos resultados; adota-se aqui apenas \SI{72,9}{\percent} (Tabela~3).

\bibitem{kumar2026}
G.~Kumar, M.~Terauds, A.~Ajayakumar~Raji, J.~Semenako, V.~Smolaninovs, P.~E.~Sics e A.~K.~Malayidinja~Poikayil~Thankappan,
\newblock Design of a lightweight edge-AI system for predictive maintenance on ESP32-S3,
\newblock \emph{Applied Sciences}, v.~16, n.~11, art.~5287, 2026.
\newblock DOI: 10.3390/app16115287.

\bibitem{landi2023}
E.~Landi, A.~Prato, A.~Fort, M.~Mugnaini, V.~Vignoli, A.~Facello, F.~Mazzoleni, M.~Murgia e A.~Schiavi,
\newblock Highly reliable multicomponent MEMS sensor for predictive maintenance management of rolling bearings,
\newblock \emph{Micromachines}, v.~14, n.~2, art.~376, 2023.
\newblock DOI: 10.3390/mi14020376.

\bibitem{liu2008}
F.~T.~Liu, K.~M.~Ting e Z.-H.~Zhou,
\newblock Isolation Forest,
\newblock \emph{Proceedings of the 2008 Eighth IEEE International Conference on Data Mining (ICDM)}, Pisa, p.~413--422, 2008.
\newblock DOI: 10.1109/ICDM.2008.17.

\bibitem{looney2017}
M.~Looney,
\newblock MEMS vibration monitoring: from acceleration to velocity,
\newblock \emph{Analog Dialogue} (Analog Devices), v.~51, n.~6, jun.\ 2017.
\newblock Literatura técnica de fabricante, não revisada por pares. O orçamento de ruído usualmente atribuído a este artigo não foi localizado em seu corpo; a atribuição permanece pendente.

\bibitem{ma2018}
S.~Ma, M.~Jiang, P.~Tao, C.~Song, J.~Wu, J.~Wang, T.~Deng e W.~Shang,
\newblock Temperature effect and thermal impact in lithium-ion batteries: a review,
\newblock \emph{Progress in Natural Science: Materials International}, v.~28, n.~6, p.~653--666, 2018.
\newblock DOI: 10.1016/j.pnsc.2018.11.002.
\newblock Texto integral não acessado (assinatura); citado a partir do resumo.

\bibitem{ma2023}
M.~Q.~Ma, Y.~Zhao, X.~Zhang e L.~Akoglu,
\newblock The need for unsupervised outlier model selection: a review and evaluation of internal evaluation strategies,
\newblock \emph{ACM SIGKDD Explorations Newsletter}, v.~25, n.~1, p.~19--35, 2023.
\newblock DOI: 10.1145/3606274.3606277.
\newblock O preprint arXiv:2104.01422 (2021) circula com título distinto e não deve ser citado como se fosse esta versão.

\bibitem{maruthi2016}
G.~S.~Maruthi e V.~Hegde,
\newblock Application of MEMS accelerometer for detection and diagnosis of multiple faults in the roller element bearings of three phase induction motor,
\newblock \emph{IEEE Sensors Journal}, v.~16, n.~1, p.~145--152, 2016.
\newblock DOI: 10.1109/JSEN.2015.2476561.

\bibitem{nabavi2025}
S.~M.~H.~Nabavi, S.~Hajforoosh e S.~M.~A.~Nabavi,
\newblock Battery selection for energy-efficient IoT devices: a comparative study of longevity across environmental conditions,
\newblock \emph{International Journal of Electrochemistry}, v.~2025, n.~1, art.~7783200, 2025.
\newblock DOI: 10.1155/ijel/7783200.

\bibitem{randall2011}
R.~B.~Randall e J.~Antoni,
\newblock Rolling element bearing diagnostics: a tutorial,
\newblock \emph{Mechanical Systems and Signal Processing}, v.~25, n.~2, p.~485--520, 2011.
\newblock DOI: 10.1016/j.ymssp.2010.07.017.

\bibitem{rodrigues2021}
J.~V.~O.~Rodrigues, M.~P.~G.~Pedroso, F.~F.~B.~Silva e R.~G.~Leão~Junior,
\newblock Performance evaluation of accelerometers ADXL345 and MPU6050 exposed to random vibrational input,
\newblock \emph{Research, Society and Development}, v.~10, n.~15, art.~e286101523082, 2021.
\newblock DOI: 10.33448/rsd-v10i15.23082.

\bibitem{scikitlearn2026}
scikit-learn developers,
\newblock \emph{sklearn.ensemble.IsolationForest} --- documentação oficial da API, versão 1.9.0,
\newblock \url{https://scikit-learn.org/stable/modules/generated/sklearn.ensemble.IsolationForest.html}. Acesso em ago.\ 2026.

\bibitem{sharma2024}
G.~Sharma, T.~Kaur, S.~K.~Mangal e A.~Kohli,
\newblock Investigating bearing and gear vibrations with a Micro-Electro-Mechanical Systems (MEMS) and machine learning approach,
\newblock \emph{Results in Engineering}, v.~24, art.~103499, 2024.
\newblock DOI: 10.1016/j.rineng.2024.103499.
\newblock Achado extraído do resumo e dos destaques depositados pelo editor; texto integral não acessado.

\bibitem{siffer2017}
A.~Siffer, P.-A.~Fouque, A.~Termier e C.~Largouët,
\newblock Anomaly detection in streams with extreme value theory,
\newblock \emph{Proceedings of the 23rd ACM SIGKDD International Conference on Knowledge Discovery and Data Mining (KDD '17)}, Halifax, p.~1067--1075, 2017.
\newblock DOI: 10.1145/3097983.3098144.

\bibitem{tractian2022}
TRACTIAN,
\newblock \emph{Smart Trac --- automated machine monitoring} --- folha de dados,
\newblock TRACTIAN, São Paulo, ago.\ 2022.
\newblock Literatura de fabricante, não revisada por pares.

\bibitem{uselis2025}
R.~\=Uselis, A.~Serackis e R.~Pomarnacki,
\newblock Signal processing optimization in resource-limited IoT for fault prediction in rotating machinery,
\newblock \emph{Electronics}, v.~14, n.~18, art.~3670, 2025.
\newblock DOI: 10.3390/electronics14183670.

\end{thebibliography}

\end{document}
